% ch1.tex -- Module 8, Chapter 1: Origins of the Elements (IQ1)
% Figures: figures/ch1_*.py (matplotlib); the rest is TikZ below.
\chapter{Origins of the Elements: the Big Bang and the Stars}\label{ch1:chap}
\covers{M8.IQ1.01, M8.IQ1.02, M8.IQ1.03, M8.IQ1.04, M8.IQ1.05, M8.IQ1.06, M8.IQ1.07}

% nuclide symbol with right-aligned mass number A and atomic number Z: \chOneNuc{A}{Z}{X}
\newlength\chOneNucW
\newlength\chOneNucV
\newcommand\chOneNuc[3]{%
  \settowidth\chOneNucW{$\scriptstyle #1$}\settowidth\chOneNucV{$\scriptstyle #2$}%
  \ifdim\chOneNucV>\chOneNucW\chOneNucW=\chOneNucV\fi
  {}^{\makebox[\chOneNucW][r]{$\scriptstyle #1$}}_{\makebox[\chOneNucW][r]{$\scriptstyle #2$}}\mathrm{#3}}

\inquiry{What evidence is there for the origins of the elements?}

\begin{goals}
  \item explain how energy turned into matter after the Big Bang, using $E = mc^2$
  \item describe Hubble's evidence for an expanding Universe and use $v = H_0 d$
  \item explain continuous, emission and absorption spectra, and use
        $\lambda_\mathrm{max} = \dfrac{b}{T}$
  \item classify stars by their spectra and read a Hertzsprung--Russell diagram
  \item describe the life of a star and the fusion reactions that make the elements
\end{goals}

%% ==================================================================
\section{From radiation to matter: the Big Bang}\label{ch1:sec:bigbang}
\covers{M8.IQ1.01, M8.IQ1.03}

\begin{bigidea}
About 13.8 billion years ago the Universe was extremely hot and dense. Since then
space has expanded and the Universe has cooled. In the hot early Universe, energy
turned into matter, as $E = mc^2$ allows.
\end{bigidea}

\begin{keywords}
  \kw{Antiparticle}{a particle with the same mass as a normal particle but the opposite charge
    (and other opposite properties).}
  \kw{Pair production}{energy (photons) turning into a particle and its antiparticle.}
  \kw{Annihilation}{a particle and its antiparticle meeting and turning into photons.}
  \kw{Nucleosynthesis}{the making of new atomic nuclei.}
\end{keywords}

The \term{Big Bang} theory says the Universe began about 13.8 billion years ago. At
first it was tiny, extremely hot and dense. It was \emph{not} an explosion into empty
space. Space itself began to stretch, everywhere at once. So there is no centre and no
edge.

As space expanded, the Universe cooled. Air rushing out of a bike tyre feels cold for
the same reason: an expanding gas cools down.

\begin{refresher}{Mass is a form of energy (Module 7)}
$E = mc^2$: mass can turn into energy and energy into mass. Because $c^2$ is huge, a
tiny mass is a lot of energy. Small energies are measured in \term{electronvolts} (eV).
One electronvolt is the energy one electron gains when it is pushed through 1 volt:
$\qty{1}{eV} = \qty{1.602e-19}{J}$. One mega-electronvolt is $\qty{1}{MeV} = \qty{e6}{eV}$.
\end{refresher}

\subsection{How energy became matter}
At first the Universe was full of \term{photons} (packets of light energy) with
enormous energies. A photon with enough energy can pass close to a nucleus, or hit
another photon. Then it can become a particle and its \term{antiparticle}. This is
\term{pair production}. For example, it can become an electron ($\mathrm{e}^-$) and a
positron ($\mathrm{e}^+$, the electron's antiparticle).

The photon's energy must be at least the energy locked in the mass of \emph{both}
particles, $mc^2$ for each one: $E \geq 2mc^2$. \Cref{ch1:fig:pairs} shows the two ways
it can happen, and the reverse process, \term{annihilation}.

\begin{figure}
  \centering
  \begin{tikzpicture}[x=1cm, y=1cm,
      photon/.style={decorate, decoration={snake, amplitude=0.8mm, segment length=2.4mm,
        post length=2mm}, line width=0.9pt, color=hscray},
      parrow/.style={-{Stealth[length=2.2mm, width=1.7mm]}, line width=1.1pt}]
    \begin{scope}
      \draw[photon, -{Stealth[length=2.2mm]}] (0,0) -- (1.55,0);
      \node[diagramlabel, above] at (0.75,0.15) {$\gamma$};
      \fill[hscmuted!55] (1.75,-0.5) circle (0.17);
      \node[font=\scriptsize, text=white] at (1.75,-0.5) {$+$};
      \node[diagramlabel, below, font=\footnotesize] at (1.75,-0.7) {nucleus};
      \draw[parrow, hscvel] (1.7,0) -- (3.0,0.75) node[right, text=hscvel] {$\mathrm{e}^-$};
      \draw[parrow, hscforce] (1.7,0) -- (3.0,-0.55) node[right, text=hscforce] {$\mathrm{e}^+$};
      \node[diagramlabel, font=\footnotesize] at (1.6,-1.6) {(a) one photon near a nucleus};
    \end{scope}
    \begin{scope}[xshift=5.0cm]
      \draw[photon, -{Stealth[length=2.2mm]}] (0,0) -- (1.3,0);
      \draw[photon, -{Stealth[length=2.2mm]}] (3.0,0) -- (1.7,0);
      \node[diagramlabel, above] at (0.6,0.15) {$\gamma$};
      \node[diagramlabel, above] at (2.4,0.15) {$\gamma$};
      \draw[parrow, hscvel] (1.5,0.1) -- (1.5,1.0) node[above, text=hscvel] {$\mathrm{e}^-$};
      \draw[parrow, hscforce] (1.5,-0.1) -- (1.5,-0.9) node[below, text=hscforce] {$\mathrm{e}^+$};
      \node[diagramlabel, font=\footnotesize] at (1.5,-1.6) {(b) two photons collide};
    \end{scope}
    \begin{scope}[xshift=10.0cm]
      \draw[parrow, hscvel] (0,0) -- (1.3,0);
      \draw[parrow, hscforce] (3.0,0) -- (1.7,0);
      \node[diagramlabel, above, text=hscvel] at (0.5,0.1) {$\mathrm{e}^-$};
      \node[diagramlabel, above, text=hscforce] at (2.5,0.1) {$\mathrm{e}^+$};
      \draw[photon, -{Stealth[length=2.2mm]}] (1.5,0.15) -- (1.5,1.15);
      \draw[photon, -{Stealth[length=2.2mm]}] (1.5,-0.15) -- (1.5,-1.0);
      \node[diagramlabel, right] at (1.65,0.85) {$\gamma$};
      \node[diagramlabel, right] at (1.65,-0.75) {$\gamma$};
      \node[diagramlabel, font=\footnotesize] at (1.5,-1.6) {(c) annihilation};
    \end{scope}
  \end{tikzpicture}
  \caption{Pair production turns energy into matter. (a) A single photon cannot make a
    pair on its own, because energy and momentum could not both stay the same. A nearby
    nucleus takes a little of the momentum. (b) No nucleus is needed. This matters
    because the very early Universe had no nuclei. (c) Two photons fly off in opposite
    directions, so the total momentum stays the same (zero).}\label{ch1:fig:pairs}
\end{figure}

While the Universe was very hot, pairs were made and destroyed all the time. As it
cooled, photons lost energy. First they could no longer make heavy particles, such as
protons and antiprotons. Within a few seconds they could not even make
electron--positron pairs. Almost all matter annihilated
with antimatter, but about one extra particle in a billion was left over. All the
matter we see today comes from this small excess. Nobody yet knows why it existed.

\begin{toolkit}{Powers of ten (scientific notation)}
\num{3.00e8} means $3.00 \times \num{100000000} = \num{300000000}$.
\begin{itemize}
  \item Multiply: add the powers, $10^a \times 10^b = 10^{a+b}$. Example:
        $\num{9.109e-31} \times \num{9.00e16} = (9.109 \times 9.00) \times 10^{-31+16}
        = \num{81.98e-15} = \num{8.198e-14}$.
  \item Divide: subtract the powers, $\dfrac{10^a}{10^b} = 10^{a-b}$.
  \item Square: square the number and double the power,
        $(\num{3.00e8})^2 = 3.00^2 \times 10^{8 \times 2} = \num{9.00e16}$.
\end{itemize}
On a calculator, type \num{9.109e-31} as \texttt{9.109 EXP -31} (or \texttt{×10\textsuperscript{x}}).
\end{toolkit}

\Cref{ch1:fig:timeline} shows the order of events. Two new words appear in it.
\term{Quarks} are tiny particles: three quarks join to make each proton or neutron
(\cref{ch5:sec:sm}). The \term{CMB} is light set free when the first atoms formed
(\cref{ch1:sec:expansion}). Before about \qty{e-43}{s}, our laws of physics cannot
describe the Universe.

The step at 3 to 20 minutes is \term{Big Bang nucleosynthesis}: the making of the first
helium nuclei. It left about \qty{75}{\percent} hydrogen and \qty{25}{\percent} helium by
mass, and a trace of lithium. Heavier elements were made later, in stars.

\begin{figure}
  \centering
  \begin{tikzpicture}[x=1mm, y=1mm,
      ev/.style={font=\footnotesize, text width=33mm, align=center, inner sep=1pt,
        text=hscink}]
    \draw[axisline, line width=0.9pt] (0,0) -- (150,0);
    \foreach \x in {6, 24.5, 43, 61.5, 80, 98.5, 117, 135.5}
      \fill[hscbrand] (\x,0) circle (1.1);
    \foreach \x in {6, 43, 80, 117} \draw[hscmuted] (\x,1.3) -- (\x,3);
    \foreach \x in {24.5, 61.5, 98.5, 135.5} \draw[hscmuted] (\x,-1.3) -- (\x,-3);
    \node[ev, anchor=south west, text width=24mm, align=left] at (0,3.2)
      {\textbf{Big Bang}\\ hot and dense};
    \node[ev, anchor=north] at (24.5,-3.2) {\textbf{\num{e-36} to \qty{e-32}{s}}\\ inflation: very fast expansion};
    \node[ev, anchor=south] at (43,3.2) {\textbf{before \qty{e-6}{s}}\\ photons make heavy particle--antiparticle pairs};
    \node[ev, anchor=north] at (61.5,-3.2) {\textbf{about \qty{e-6}{s}}\\ quarks join into protons and neutrons};
    \node[ev, anchor=south] at (80,3.2) {\textbf{3 to 20 minutes} (\qty{e9}{K})\\ protons and neutrons join into helium nuclei};
    \node[ev, anchor=north] at (98.5,-3.2) {\textbf{\num{380000} years} (\qty{3000}{K})\\ atoms form; light (CMB) set free};
    \node[ev, anchor=south] at (117,3.2) {\textbf{200 million years}\\ first stars and galaxies};
    \node[ev, anchor=north east, text width=26mm, align=right] at (150,-3.2)
      {\textbf{13.8 billion years}\\ today (\qty{2.7}{K})};
  \end{tikzpicture}
  \caption{A timeline of the Universe (not to scale). The temperature falls all the
    time. Electron--positron pairs were still made until a few seconds after the Big
    Bang.}\label{ch1:fig:timeline}
\end{figure}

\begin{toolkit}{Rearranging a formula}
To make another symbol the subject, do the same thing to both sides until that symbol
is alone.
\begin{itemize}
  \item Make $m$ the subject of $E = mc^2$: divide both sides by $c^2$. Then
        $\dfrac{E}{c^2} = m$, that is $m = \dfrac{E}{c^2}$.
  \item When the symbol is on the bottom, first move it to the top. Make $T$ the subject
        of $\lambda_\mathrm{max} = \dfrac{b}{T}$ (Wien's law, \cref{ch1:sec:spectra}):
        multiply both sides by $T$ ($\lambda_\mathrm{max}T = b$). Then divide both sides by
        $\lambda_\mathrm{max}$: $T = \dfrac{b}{\lambda_\mathrm{max}}$.
\end{itemize}
\end{toolkit}

\begin{formulacard}{Mass--energy equivalence}
  \cardformula{E = mc^2}
  \begin{symbols}
    \sym{$E$}{energy}{\si{J}}
    \sym{$m$}{mass (made, or turned into energy)}{\si{kg}}
    \sym{$c$}{speed of light, \qty{3.00e8}{m.s^{-1}}}{\si{m.s^{-1}}}
  \end{symbols}
  \inwords{mass and energy are two forms of the same thing.}
  \whentouse{energy turning into matter (pair production) or matter into energy
    (annihilation; fusion, where light nuclei join).}
  \whennottouse{kinetic energy of everyday objects (use $K = \frac{1}{2}mv^2$).}
  \rearranged{$m = \dfrac{E}{c^2}$ \qquad pair: $E_\mathrm{min} = 2mc^2$}
  \tensecond{$m = \qty{1.0e-3}{kg}$: $E = \num{1.0e-3} \times (\num{3.00e8})^2 = \qty{9.0e13}{J}$.}
  \trap{$m$ in kg gives $E$ in J; divide by \qty{1.602e-19}{J} for eV. For a pair use $2m$.}
\end{formulacard}

\begin{workedexample}{Making a proton and an antiproton}\label{ch1:we:proton}
A gamma photon near a heavy nucleus makes a proton--antiproton pair. Ignore the small
energy the nucleus takes. Find the smallest photon energy, in J and in MeV.
\given{$m_\mathrm{p} = \qty{1.673e-27}{kg}$, \ $c = \qty{3.00e8}{m.s^{-1}}$, \
  $\qty{1}{eV} = \qty{1.602e-19}{J}$}
\find{the minimum photon energy $E$}
\formula{$E = mc^2$ with $m = 2m_\mathrm{p}$}
\substitute{\eqlines{E &= 2 \times \qty{1.673e-27}{kg} \times (\qty{3.00e8}{m.s^{-1}})^2
  = \qty{3.011e-10}{J}}\par
  \eqlines{E &= \frac{\qty{3.011e-10}{J}}{\qty{1.602e-19}{J.eV^{-1}}} = \qty{1.880e9}{eV}
  \quad\text{(change J to eV)}\\
  E &= \frac{\qty{1.880e9}{eV}}{\num{e6}} = \qty{1.880e3}{MeV}
  \quad\text{($\qty{1}{MeV} = \qty{e6}{eV}$)}}}
\answer{\final{$E = \qty{3.01e-10}{J} = \qty{1.88e3}{MeV}$}}
\checkstep{rough size, with $c^2 \approx 10^{17}$:
  $2 \times \num{1.7e-27} \times 10^{17} \approx \qty{3.4e-10}{J}$, close \cmark.}
\end{workedexample}

\begin{refresher}{Photons, frequency and wavelength (Module 7)}
Light comes in tiny packets of energy called photons. One photon has energy $E = hf$.
Here $h = \qty{6.626e-34}{J.s}$ (Planck's constant) and $f$ is the frequency in hertz
(\si{Hz}). For all waves $v = f\lambda$. For light $v = c$, so $c = f\lambda$ and
$\lambda = \dfrac{c}{f}$. A higher frequency means a shorter wavelength and more energy
in each photon. Gamma rays have the most energy per photon; radio waves have the least.
\end{refresher}

\begin{workedexample}{Two photons make an electron and a positron}\label{ch1:we:twophotons}
Two identical photons hit each other head-on. Together they have the smallest energy
that can make an electron and a positron. Find the frequency and wavelength of each
photon. Which part of the spectrum is it?
\given{$m_\mathrm{e} = \qty{9.109e-31}{kg}$, \ $h = \qty{6.626e-34}{J.s}$, \
  $c = \qty{3.00e8}{m.s^{-1}}$}
\find{$f$ and $\lambda$ of each photon}
\formula{the pair needs $2m_\mathrm{e}c^2$, so each photon gives half: $E = m_\mathrm{e}c^2$;
  \ then $E = hf$ \ and \ $v = f\lambda$ (with $v = c$)}
\substitute{\eqlines{E &= \tfrac{1}{2} \times 2m_\mathrm{e}c^2 = m_\mathrm{e}c^2
  \quad\text{(the two photons share the energy equally)}\\
  E &= \qty{9.109e-31}{kg} \times (\qty{3.00e8}{m.s^{-1}})^2 = \qty{8.198e-14}{J}}\par
  \eqlines{f &= \frac{E}{h} = \frac{\qty{8.198e-14}{J}}{\qty{6.626e-34}{J.s}} = \qty{1.2372e20}{Hz}\\
  \lambda &= \frac{c}{f} = \frac{\qty{3.00e8}{m.s^{-1}}}{\qty{1.2372e20}{Hz}} = \qty{2.4248e-12}{m}\\
  &\quad\text{(keep extra figures until the last step)}}}
\answer{\final{$f = \qty{1.24e20}{Hz}$, \ $\lambda = \qty{2.42e-12}{m}$}: gamma rays.}
\checkstep{gamma rays are the most energetic photons \cmark.}
\end{workedexample}

\begin{mistakes}
\begin{itemize}
  \item \xmark\ ``The Big Bang was an explosion from a centre.'' \cmark\ Space itself
        expanded everywhere.
  \item \xmark\ Using $m$ instead of $2m$ for a pair. \cmark\ $E_\mathrm{min} = 2mc^2$.
  \item \xmark\ ``Annihilation makes one photon.'' \cmark\ At least two (momentum).
\end{itemize}
\end{mistakes}

\begin{examtip}
``Describe the processes that turned radiation into matter after the Big Bang''
(4--5 marks). Full marks need these linked points:
\begin{enumerate}
  \item $E = mc^2$: energy can become mass.
  \item Pair production: a photon (or two photons) with $E \geq 2mc^2$ makes a particle
        and its antiparticle.
  \item As space expanded, the Universe cooled. Photons became too weak to make pairs.
  \item Matter and antimatter annihilated, but a tiny excess of matter was left.
  \item The order: quarks $\to$ protons and neutrons (about \qty{e-6}{s}) $\to$ helium
        nuclei (3--20 minutes) $\to$ atoms (\num{380000} years).
\end{enumerate}
\end{examtip}

\begin{quickcheck}
  \item Put in the order they formed: atoms, first stars, helium nuclei, protons, quarks.
  \item Find the minimum energy, in MeV, of one photon that makes an electron and a
        positron near a heavy nucleus. Ignore the energy the nucleus takes.
  \item Explain why, as the Universe cooled, pair production of protons stopped before
        pair production of electrons.
\end{quickcheck}

%% ==================================================================
\section{Evidence for an expanding Universe}\label{ch1:sec:expansion}
\covers{M8.IQ1.02, M8.IQ1.01}

\begin{bigidea}
Light from distant galaxies is shifted to longer, redder wavelengths. The further away
a galaxy is, the faster it moves away. So the Universe is expanding, and long ago it
was small, hot and dense.
\end{bigidea}

\begin{keywords}
  \kw{Galaxy}{a huge group of billions of stars held together by gravity.}
  \kw{Redshift}{spectral lines moved to longer wavelengths, because the source moves away
    or because space stretches while the light travels.}
  \kw{Recession velocity ($v$)}{the speed at which a galaxy moves away from us.}
  \kw{Megaparsec (Mpc)}{a distance unit: $\qty{1}{Mpc} = \qty{3.086e22}{m}$ (3.26 million light-years).}
\end{keywords}

\subsection{Redshift}
Light from a star or a galaxy can be spread out into a rainbow of colours, called a
\term{spectrum}. Gas in the star absorbs a few exact wavelengths, so the rainbow has
thin dark lines in it. These \term{spectral lines} are at fixed wavelengths for each
element, so they work like a fingerprint. (The top strip of \cref{ch1:fig:redshift}
shows hydrogen's lines. \Cref{ch1:sec:spectra} explains why they form.)

\begin{refresher}{The Doppler effect (Year 11)}
A siren sounds lower in pitch (a deeper note) as an ambulance drives away. The
ambulance moves away a little between one wave and the next, so the waves are
stretched. They have a longer wavelength and a lower frequency. Light from a moving
source behaves in the same way.
\end{refresher}

When a galaxy moves away from us, its light is stretched to longer wavelengths. In
visible light, red has the longest wavelength, so we call this \term{redshift}. Light
from a source moving towards us is squashed to shorter wavelengths: \term{blueshift}.

In \cref{ch1:fig:redshift} every line has moved. But the pattern of lines is the same,
so we can still tell which line is which. A bigger shift means a faster speed.

\begin{figure}
  \centering
  \includegraphics{ch1_redshift.pdf}
  \caption{Redshift. Every line moves by the same \emph{fraction} of its wavelength
    (here \qty{3.5}{\percent}).}\label{ch1:fig:redshift}
\end{figure}

\begin{formulacard}[notonsheet]{Redshift (low speeds)}
  \cardformula{z = \frac{\Delta\lambda}{\lambda_0} = \frac{\lambda - \lambda_0}{\lambda_0}
    \approx \frac{v}{c}}
  \begin{symbols}
    \sym{$z$}{redshift}{no unit}
    \sym{$\lambda_0$, $\lambda$}{wavelength in the laboratory; wavelength observed}{\si{m} (or \si{nm})}
    \sym{$\Delta\lambda$}{shift in wavelength, $\lambda - \lambda_0$}{\si{m} (or \si{nm})}
    \sym{$v$}{recession velocity}{\si{m.s^{-1}}}
    \sym{$c$}{speed of light, \qty{3.00e8}{m.s^{-1}}}{\si{m.s^{-1}}}
  \end{symbols}
  \inwords{the fraction a wavelength stretches = the galaxy's speed as a fraction of $c$.}
  \whentouse{a galaxy moving much slower than light ($v$ below about $0.1c$).}
  \whennottouse{sound (use the Doppler formula on the sheet).}
  \rearranged{$v = \dfrac{c\,\Delta\lambda}{\lambda_0}$ \qquad
    $\Delta\lambda = \dfrac{\lambda_0 v}{c}$}
  \tensecond{$\lambda_0 = \qty{500.0}{nm}$, $\lambda = \qty{505.0}{nm}$: $z = 0.010$,
    so $v = 0.010c = \qty{3.0e6}{m.s^{-1}}$.}
  \derive{the source sends one wave every period $T$, so $\lambda_0 = cT$. Moving away at
    $v$, it is $vT$ further away when it sends the next wave, so each wave is stretched by
    $\Delta\lambda = vT$. Then $\dfrac{\Delta\lambda}{\lambda_0} = \dfrac{vT}{cT} = \dfrac{v}{c}$.}
  \trap{$\lambda$ and $\lambda_0$ in the same unit. $v$ comes out in \si{m.s^{-1}}:
    divide by 1000 for \si{km.s^{-1}}.}
\end{formulacard}

\subsection{Hubble's evidence}
\begin{itemize}
  \item \emph{Speeds}: from about 1912, Vesto Slipher measured the spectra of faint,
        fuzzy patches called ``spiral nebulae'' (we now know they are galaxies). Most
        were redshifted, so most were moving away.
  \item \emph{Distances}: a \term{Cepheid variable} is a star that gets brighter and
        dimmer again and again, in a regular cycle. Henrietta Leavitt (1912) found that
        the time for one cycle tells us the star's luminosity (true brightness).
        Comparing this with how bright it looks from Earth gives its distance.
  \item In 1924 Edwin Hubble used Cepheids to measure the distance to the Andromeda
        ``nebula''. It is far outside the Milky Way, so it is another galaxy.
  \item In 1929 Hubble plotted speed against distance. Speed was proportional to
        distance: a galaxy twice as far away moved away twice as fast
        (\cref{ch1:fig:hubble}).
\end{itemize}

\begin{toolkit}{Gradient of a straight-line graph}
The gradient (slope) of a straight line $= \dfrac{\text{rise}}{\text{run}}
= \dfrac{\text{change in the vertical quantity}}{\text{change in the horizontal quantity}}$.
\begin{enumerate}
  \item Pick two points on the line, far apart. Use points on the line, not data points.
  \item Read both coordinates, with units.
  \item Divide the rise by the run. The unit is (vertical unit) $\div$ (horizontal unit).
\end{enumerate}
\Cref{ch1:fig:hubble}: from (\qty{40}{Mpc}, \qty{2796}{km.s^{-1}}) to
(\qty{120}{Mpc}, \qty{8388}{km.s^{-1}}), the gradient is
$\dfrac{(8388 - 2796)\ \si{km.s^{-1}}}{(120 - 40)\ \si{Mpc}} = \dfrac{5592}{80}
= \qty{69.9}{km.s^{-1}.Mpc^{-1}}$.
\end{toolkit}

\begin{sidebyside}[0.54]
  \centering
  \begin{tikzpicture}
    \begin{axis}[hsc, width=0.82\linewidth, height=0.64\linewidth,
        xmin=0, xmax=150, ymin=0, ymax=11000,
        xlabel={Distance $d$ (\si{Mpc})}, ylabel={$v$ (\si{km.s^{-1}})},
        ytick={0,2000,...,10000},
        yticklabel style={/pgf/number format/.cd, 1000 sep={\,}}]
      \addplot[hscseriesB, dashed, line width=1pt, domain=0:150, samples=2] {69.9*x};
      \addplot[hscseriesA, only marks, mark=*, mark size=2pt]
        coordinates {(22,1600) (48,3200) (75,5390) (104,7040) (140,9920)};
      \draw[hscmuted, line width=0.6pt] (axis cs:40,2796) -- (axis cs:120,2796)
        -- (axis cs:120,8388);
      \node[font=\scriptsize, below] at (axis cs:80,2796) {$\Delta d = \qty{80}{Mpc}$};
      \node[font=\scriptsize, align=left, anchor=west] at (axis cs:121,5300)
        {$\Delta v =$\\\qty{5590}{km.s^{-1}}};
    \end{axis}
  \end{tikzpicture}
  \captionof{figure}{Recession velocity against distance for five galaxies (data in
    \cref{ch1:sec:experiments}).}\label{ch1:fig:hubble}
\nextside
  \begin{itemize}
    \item The points lie near a straight line through the origin. So $v \propto d$,
          which means ``$v$ is proportional to $d$'': if $d$ doubles, $v$ doubles.
    \item The gradient is the \term{Hubble constant}, $H_0$:
          $\dfrac{\qty{5590}{km.s^{-1}}}{\qty{80}{Mpc}} = \qty{69.9}{km.s^{-1}.Mpc^{-1}}$.
    \item If the axes are swapped ($d$ up, $v$ across), the gradient is
          $\dfrac{d}{v} = \dfrac{1}{H_0}$.
  \end{itemize}
\end{sidebyside}

\subsection{What the expansion means}
All distant galaxies move away from us. That does not put Earth at the centre. The
\term{cosmological principle} says that, on a very large scale, the Universe looks the
same from every place and in every direction. From any galaxy, all other galaxies would
be seen moving away.

Picture a loaf of raisin bread rising in a warm oven. The dough between the raisins
grows, so the raisins spread apart. Stand on any raisin: a raisin \qty{2}{cm} away
moves off twice as fast as one \qty{1}{cm} away. Every raisin sees the same thing, so
none of them is the centre. In the same way, \emph{space itself} is stretching.

Run the expansion backwards: everything was once together, hot and dense. The Universe
had a beginning.

\begin{formulacard}[notonsheet]{Hubble's law and the age of the Universe}
  \cardformula{v = H_0 d \qquad\qquad t \approx \frac{1}{H_0}}
  \begin{symbols}
    \sym{$v$}{recession velocity}{\si{m.s^{-1}} (usually given in \si{km.s^{-1}})}
    \sym{$d$}{distance to the galaxy}{\si{m} (usually given in \si{Mpc})}
    \sym{$H_0$}{Hubble constant (about 70 \si{km.s^{-1}.Mpc^{-1}}; given in the question)}{\si{s^{-1}} (usually given in \si{km.s^{-1}.Mpc^{-1}})}
    \sym{$t$}{age of the Universe (use $H_0$ in \si{s^{-1}})}{\si{s}}
  \end{symbols}
  \inwords{a galaxy twice as far away moves away twice as fast.}
  \whentouse{distant galaxies; $t \approx \frac{1}{H_0}$ for a rough age.}
  \whennottouse{stars in our galaxy, or very close galaxies (Andromeda moves towards us).}
  \rearranged{$d = \dfrac{v}{H_0}$ \qquad $H_0 = \dfrac{v}{d}$ (gradient of a $v$--$d$ graph)}
  \tensecond{$d = \qty{40}{Mpc}$, $H_0 = 70$: $v = 70 \times 40 = \qty{2800}{km.s^{-1}}$.}
  \derive{$v = H_0 d$ comes from the straight-line graph. For the age, assume each
    galaxy has always moved at the same speed:
    \begin{align*}
      v &= \frac{d}{t} && \text{speed = distance $\div$ time}\\
      H_0 d &= \frac{d}{t} && \text{replace $v$ with $H_0 d$}\\
      H_0 &= \frac{1}{t} && \text{divide both sides by $d$}\\
      t &= \frac{1}{H_0} && \text{flip both sides}
    \end{align*}}
  \trap{never mix \si{m} and \si{km}. Change $H_0$ into \si{s^{-1}} before using
    $t \approx \frac{1}{H_0}$.}
\end{formulacard}

\subsection{More evidence: the CMB and the light elements}
In 1965 Arno Penzias and Robert Wilson found a faint microwave signal. It came equally
from every direction. It is the \term{cosmic microwave background} (CMB): light left
over from the early Universe.

Before \num{380000} years, the Universe was full of free electrons. They scattered light
all the time, like fog scatters a car's headlights. Then the Universe cooled to about
\qty{3000}{K}. Electrons joined nuclei to make atoms, the ``fog'' cleared, and light
could travel freely. That light is the CMB. Then it was mostly infrared and visible light.

Since then, the expansion of space has stretched its wavelengths about 1100 times, into
microwaves. Today it matches a black body at \qty{2.7}{K}. (A black body is an ideal
glowing object whose light depends only on its temperature: see \cref{ch1:sec:spectra}
and \cref{ch1:we:cmb}.) The Big Bang theory predicted this light in 1948, before it was
found.

Big Bang nucleosynthesis predicts \qty{75}{\percent} hydrogen and \qty{25}{\percent}
helium by mass. The oldest stars and gas clouds have almost exactly these amounts. Stars
could not have made that much helium in 13.8 billion years.

\begin{workedexample}{A galaxy's speed and its redshift}\label{ch1:we:galaxy}
A galaxy is \qty{150}{Mpc} away. Take $H_0 = \qty{70.0}{km.s^{-1}.Mpc^{-1}}$. (a) Find
its recession velocity. (b) Hydrogen's H-alpha line is at \qty{656.3}{nm} in the
laboratory. At what wavelength do we observe it from this galaxy?
\given{$d = \qty{150}{Mpc}$, \ $H_0 = \qty{70.0}{km.s^{-1}.Mpc^{-1}}$, \
  $\lambda_0 = \qty{656.3}{nm}$, \ $c = \qty{3.00e8}{m.s^{-1}}$}
\find{(a) $v$; \ (b) the observed wavelength $\lambda$}
\formula{$v = H_0 d$, \ then $\Delta\lambda = \dfrac{\lambda_0 v}{c}$}
\substitute{(a)\ \eqlines{v &= \qty{70.0}{km.s^{-1}.Mpc^{-1}} \times \qty{150}{Mpc}
  = \qty{10500}{km.s^{-1}}\\
  &\quad\text{(the Mpc cancel)}\\
  v &= \num{10500} \times \qty{1000}{m.s^{-1}} = \qty{1.05e7}{m.s^{-1}}\\
  &\quad\text{($\times\,1000$ changes km to m)}}\par
  (b)\ \eqlines{\Delta\lambda &= \qty{656.3}{nm} \times \frac{\qty{1.05e7}{m.s^{-1}}}{\qty{3.00e8}{m.s^{-1}}}
  = \qty{22.97}{nm}\\
  \lambda &= \lambda_0 + \Delta\lambda \quad\text{(add the shift)}\\
  &= \qty{656.3}{nm} + \qty{22.97}{nm} = \qty{679.27}{nm}}}
\answer{(a) \final{$v = \qty{1.05e4}{km.s^{-1}}$} \quad (b) \final{$\lambda = \qty{679.3}{nm}$}}
\checkstep{$v$ is only \qty{3.5}{\percent} of $c$, so the low-speed formula is fine.
  This is the galaxy in \cref{ch1:fig:redshift} \cmark.}
\end{workedexample}

\begin{toolkit}{Changing $H_0$ into \si{s^{-1}}}
$\si{Mpc^{-1}}$ means ``per megaparsec''. $\qty{1}{Mpc} = \qty{3.086e22}{m}
= \qty{3.086e19}{km}$ (divide by 1000 to change m to km). So $H_0 = 70$ means
\qty{70}{km.s^{-1}} for every \qty{3.086e19}{km}:
\[
  H_0 = \frac{\qty{70}{km.s^{-1}}}{\qty{3.086e19}{km}} = \qty{2.27e-18}{s^{-1}}
\]
The km on the top and the bottom cancel, leaving \si{s^{-1}}.
Years to seconds: $\qty{1}{year} = 365.25 \text{ days} \times 24 \text{ hours} \times
\qty{3600}{s} = \qty{3.156e7}{s}$.
\end{toolkit}

\begin{workedexample}{How old is the Universe?}\label{ch1:we:age}
One team measured $H_0 = \qty{73}{km.s^{-1}.Mpc^{-1}}$. Estimate the age of the
Universe in years, and state your assumption.
\given{$H_0 = \qty{73}{km.s^{-1}.Mpc^{-1}}$, \ $\qty{1}{Mpc} = \qty{3.086e19}{km}$, \
  $\qty{1}{year} = \qty{3.156e7}{s}$}
\find{the age $t$ in years}
\formula{$t \approx \dfrac{1}{H_0}$, with $H_0$ in \si{s^{-1}}}
\substitute{\eqlines{H_0 &= \frac{\qty{73}{km.s^{-1}}}{\qty{3.086e19}{km}}
  = \qty{2.366e-18}{s^{-1}} \quad\text{(change units first)}}\par
  \eqlines{t &= \frac{1}{\qty{2.366e-18}{s^{-1}}} = \qty{4.227e17}{s}\\
  t &= \frac{\qty{4.227e17}{s}}{\qty{3.156e7}{s.year^{-1}}} = \qty{1.339e10}{years}
  \quad\text{(seconds to years)}}}
\answer{\final{$t \approx \qty{1.3e10}{years}$}, assuming the galaxies have always moved
  at the same speed.}
\checkstep{close to the accepted 13.8 billion years \cmark. A larger $H_0$ gives a
  younger Universe.}
\end{workedexample}

\begin{mistakes}
\begin{itemize}
  \item \xmark\ ``Hubble saw all the stars moving away.'' \cmark\ He measured
        \emph{galaxies}.
  \item \xmark\ ``Everything moves away from us, so we are at the centre.'' \cmark\ Every
        galaxy sees the same, because space expands everywhere.
  \item \xmark\ Using $H_0$ in \si{km.s^{-1}.Mpc^{-1}} in $t = \frac{1}{H_0}$.
        \cmark\ Change it into \si{s^{-1}} first.
\end{itemize}
\end{mistakes}

\begin{examtip}
``Outline the evidence that led Hubble to conclude that the Universe is expanding''
(3--4 marks). Write: redshift of galaxy spectra (they move away; bigger shift = faster);
distances from Cepheid variables; $v \propto d$, so $v = H_0 d$; with the cosmological
principle, space itself expands. Write ``galaxies'', never ``stars''. The CMB and the
helium abundance are \emph{extra} Big Bang evidence, not Hubble's.
\end{examtip}

\begin{quickcheck}
  \item A calcium line is at \qty{393.4}{nm} in the laboratory and at \qty{397.3}{nm} in a
        galaxy's spectrum. Find the galaxy's speed and distance ($H_0 = 70$
        \si{km.s^{-1}.Mpc^{-1}}).
  \item An astronomer lives on a planet in a galaxy \qty{300}{Mpc} from us. She measures
        the redshifts of many galaxies, including our Milky Way. Predict what she finds
        ($H_0 = 70$ \si{km.s^{-1}.Mpc^{-1}}). What does this tell us about the
        ``centre'' of the Universe?
  \item A classmate says: ``The CMB is microwave radiation, so it must come from
        something cold, like dust between the stars.'' The CMB comes equally from every
        direction, and today its peak wavelength is about \qty{1}{mm}. Explain where the
        CMB really came from, and why it is now microwave radiation.
\end{quickcheck}

%% ==================================================================
\section{Spectra and the classification of stars}\label{ch1:sec:spectra}
\covers{M8.IQ1.04, M8.IQ1.05}

\begin{bigidea}
A hot, dense object gives a continuous spectrum that depends only on its temperature.
Gases give bright or dark lines that identify the element. The dark lines in starlight
sort stars by temperature into the classes O B A F G K M.
\end{bigidea}

\begin{keywords}
  \kw{Spectrum}{light spread out into its different wavelengths (colours).}
  \kw{Black body}{an ideal object that absorbs all light; its glow depends only on its temperature.}
  \kw{Energy level}{one of the few energies an electron in an atom is allowed to have.}
  \kw{Spectral class}{a letter (O, B, A, F, G, K, M) for stars with similar spectra.}
\end{keywords}

\subsection{Three kinds of spectrum}
A prism or a diffraction grating spreads light into a spectrum
(\cref{ch1:fig:spectra}). A \term{continuous spectrum} has every wavelength. An
\term{emission spectrum} has bright lines only. An \term{absorption spectrum} is a
continuous spectrum with dark lines. Each element has its own set of lines, so the
lines tell us which elements are there.

\begin{figure}
  \centering
  \includegraphics{ch1_spectra.pdf}
  \caption{The three kinds of spectrum. Hydrogen's dark lines are at exactly the same
    wavelengths as its bright lines.}\label{ch1:fig:spectra}
\end{figure}

Line spectra come from electrons in atoms. Electrons can only be in certain
\term{energy levels} (\cref{ch1:fig:levels}; more in \cref{ch3:sec:bohr}).

\begin{sidebyside}[0.42]
  \centering
  \begin{tikzpicture}[x=0.8cm, y=1cm,
      photon/.style={decorate, decoration={snake, amplitude=0.7mm, segment length=2.2mm,
        post length=1.8mm}, line width=0.8pt, color=hscray, -{Stealth[length=2mm]}}]
    \foreach \y/\n in {0/1, 1.5/2, 2.25/3}{
      \draw[hscink, line width=0.8pt] (0,\y) -- (5.2,\y);
      \node[diagramlabel, right] at (5.2,\y) {$n = \n$};}
    \draw[axisline] (-0.35,-0.1) -- (-0.35,2.6) node[above, font=\footnotesize] {energy};
    \draw[-{Stealth[length=2mm]}, line width=1pt, hscvel] (0.9,2.25) -- (0.9,1.5);
    \fill[hscvel] (0.9,2.25) circle (0.08);
    \draw[photon] (1.05,1.87) -- (2.15,1.87);
    \node[diagramlabel, font=\footnotesize] at (1.3,0.75) {emission};
    \draw[photon] (2.95,1.87) -- (4.0,1.87);
    \draw[-{Stealth[length=2mm]}, line width=1pt, hscvel] (4.25,1.5) -- (4.25,2.25);
    \fill[hscvel] (4.25,1.5) circle (0.08);
    \node[diagramlabel, font=\footnotesize] at (3.75,0.75) {absorption};
  \end{tikzpicture}
  \captionof{figure}{An electron jumps between energy levels and gives out or takes in
    one photon.}\label{ch1:fig:levels}
\nextside
  \begin{itemize}
    \item \emph{Excitation}: in a hot gas or a discharge tube, collisions or an electric
          current lift electrons to higher levels.
    \item \emph{Emission}: an electron drops to a lower level and gives out one photon,
          $E = hf$ = the gap.
    \item \emph{Absorption}: a photon with exactly the gap energy lifts an electron up.
          The energy is given out again in random directions, so the line looks dark.
    \item Same gaps, so the same wavelengths.
  \end{itemize}
\end{sidebyside}

A \term{black body} absorbs all light that falls on it, and glows when hot. Its spectrum
depends only on its temperature. In a hot, dense object the particles are packed close
together and jostle all the time. So light of every wavelength is given out, not only
the few wavelengths that match the gaps between energy levels. Stars are very nearly
black bodies. In
\cref{ch1:fig:blackbody}, a hotter body gives out more light at every wavelength. Its
peak also moves to shorter wavelengths, so heated metal glows red, then orange, then white.
\Cref{ch1:tab:spectra} compares the three kinds.

\begin{figure}
  \centering
  \includegraphics{ch1_blackbody.pdf}
  \caption{Black body spectra for three temperatures.}\label{ch1:fig:blackbody}
\end{figure}

\begin{table}
  \caption{Comparing the three kinds of spectrum.}\label{ch1:tab:spectra}
  \centering\small
  \begin{tabular}{@{}L{24mm}L{36mm}L{36mm}L{37mm}@{}}
    \toprule
    \colhead{Spectrum} & \colhead{Made by} & \colhead{Looks like} & \colhead{Depends on}\\
    \midrule
    Continuous (black body) & a hot, dense object (filament, inside of a star) & every
      colour; a smooth curve with one peak & temperature only\\
    Emission & a hot gas at low pressure (neon sign) & bright lines on a dark background &
      which element (its energy levels)\\
    Absorption & cooler gas in front of a hot source & a rainbow with dark lines & which
      element; same places as its bright lines\\
    \bottomrule
  \end{tabular}
\end{table}

\begin{toolkit}{Nanometres and metres}
$\qty{1}{nm} = \qty{e-9}{m}$.
\begin{itemize}
  \item nm to m: multiply by $10^{-9}$. Example:
        $\qty{362}{nm} = 362 \times \qty{e-9}{m} = \qty{3.62e-7}{m}$.
  \item m to nm: multiply by $10^{9}$. Example:
        $\qty{5.00e-7}{m} = \num{5.00e-7} \times 10^9\ \si{nm} = \qty{500}{nm}$.
\end{itemize}
\end{toolkit}

\begin{formulacard}{Wien's displacement law}
  \cardformula{\lambda_\mathrm{max} = \frac{b}{T}}
  \begin{symbols}
    \sym{$\lambda_\mathrm{max}$}{wavelength where the spectrum peaks}{\si{m}}
    \sym{$b$}{Wien's constant, \qty{2.898e-3}{m.K} (data sheet)}{\si{m.K}}
    \sym{$T$}{surface temperature}{\si{K}}
  \end{symbols}
  \inwords{the hotter the object, the shorter its peak wavelength.}
  \whentouse{black bodies: stars, the CMB, hot filaments.}
  \whennottouse{emission or absorption lines.}
  \rearranged{$T = \dfrac{b}{\lambda_\mathrm{max}}$}
  \tensecond{an orange star, $T = \qty{4500}{K}$: $\lambda_\mathrm{max} =
    \dfrac{\num{2.898e-3}}{4500} = \qty{6.44e-7}{m} = \qty{644}{nm}$.}
  \trap{$T$ in kelvin, not \si{\degreeCelsius}. $\lambda_\mathrm{max}$ comes out in metres
    ($\times 10^9$ for nm).}
\end{formulacard}

\subsection{Classifying stars by their spectra}
Starlight gives an absorption spectrum. The hot, dense inside of the star makes a
continuous spectrum. Cooler gas in the star's outer layers then absorbs some
wavelengths. Astronomers sort stars into \term{spectral classes} by these lines
(\cref{ch1:tab:classes}). Hottest to coolest: O B A F G K M (``\textbf{O}h \textbf{B}e
\textbf{A} \textbf{F}ine \textbf{G}uy, \textbf{K}iss \textbf{M}e'').

Three words in the table: a \term{neutral} atom still has all its electrons. An
\term{ionised} atom has lost one or more electrons. Astronomers call every element
heavier than helium a ``metal''.

\begin{table}
  \caption{Spectral classes (temperatures are approximate). $\lambda_\mathrm{max}$ is from
    Wien's law.}\label{ch1:tab:classes}
  \centering\small
  \begin{tabular}{@{}C{9mm}L{27mm}L{20mm}L{17mm}L{58mm}@{}}
    \toprule
    \colhead{Class} & \colhead{Temperature (K)} & \colhead{Colour} &
      \colhead{$\lambda_\mathrm{max}$ (nm)} & \colhead{Main absorption lines}\\
    \midrule
    O & above \num{30000} & blue & below 97 & ionised helium; weak hydrogen\\
    B & \num{10000}--\num{30000} & blue-white & 97--290 & neutral helium; stronger hydrogen\\
    A & \num{7500}--\num{10000} & white & 290--386 & hydrogen strongest\\
    F & \num{6000}--\num{7500} & yellow-white & 386--483 & weaker hydrogen; ionised metals\\
    G & \num{5000}--\num{6000} & yellow & 483--580 & ionised calcium; many metals (Sun)\\
    K & \num{3500}--\num{5000} & orange & 580--828 & neutral metals; some molecules\\
    M & below \num{3500} & red & above 828 & molecules (titanium oxide); very weak hydrogen\\
    \bottomrule
  \end{tabular}
\end{table}

Annie Jump Cannon set this order (about 1901--1912). In 1925 Cecilia Payne showed that
the classes differ because of \emph{temperature}: most stars are made of almost the same
mix, mainly hydrogen and helium. Each class has \term{subclasses} from 0 (hottest) to 9
(coolest); the Sun is G2.

Why are the hydrogen Balmer lines (hydrogen's visible lines) strongest in class A?
Visible hydrogen lines come only from electrons that start in level $n = 2$.
\begin{itemize}
  \item Cool stars (K, M): nearly all electrons sit in $n = 1$, so the lines are weak.
  \item About \qty{10000}{K} (A): the largest share of electrons is in $n = 2$, so the
        lines are strongest.
  \item Hot stars (O): hydrogen is ionised (its electron is knocked off), so the
        lines are weak.
\end{itemize}

Spectra tell us more than the elements:
\begin{itemize}
  \item \emph{Line positions}: which elements are there.
  \item \emph{Which lines appear, and how strong (dark) they are}: the surface
        temperature. (An element's lines do not move when the temperature changes.)
  \item \emph{A shift of all the lines}: motion towards us (blueshift) or away (redshift).
  \item \emph{Line width}: how dense the star's outer gas is. Broad lines mean dense gas,
        or a star that spins fast. Narrow lines mean thin gas.
\end{itemize}
A small, dense main sequence star has broad lines. A huge giant star has thin outer gas
and narrow lines. (Giants and the main sequence are explained in \cref{ch1:sec:hr}.) So
line width gives the star's \term{luminosity class}, a Roman number for its size group:
I = supergiant, III = giant, V = main sequence. The Sun is G2V.

\begin{workedexample}{Classifying a star}\label{ch1:we:classify}
A star's continuous spectrum peaks at \qty{362}{nm}. Find its surface temperature,
spectral class and colour.
\given{$\lambda_\mathrm{max} = \qty{362}{nm} = \qty{3.62e-7}{m}$, \ $b = \qty{2.898e-3}{m.K}$}
\find{the surface temperature $T$, then the class}
\formula{$T = \dfrac{b}{\lambda_\mathrm{max}}$}
\substitute{$T = \dfrac{\qty{2.898e-3}{m.K}}{\qty{3.62e-7}{m}} = \qty{8005.5}{K}$}
\answer{\final{$T = \qty{8.01e3}{K}$}: class A (\num{7500} to \qty{10000}{K}), white.}
\checkstep{hotter than the Sun, and its peak is shorter than the Sun's \qty{500}{nm}
  \cmark. Expect strong hydrogen lines.}
\end{workedexample}

% Added by the coverage audit (M8.IQ1.04, M8.IQ1.05): numbers in
\begin{workedexample}{Reading the lines in a star's spectrum}\label{ch1:we:lines}
A star's absorption spectrum has dark lines at 393.4, 396.8, 486.1, 589.0, 589.6 and
\qty{656.3}{nm}. The lines at 393.4 and \qty{396.8}{nm} are the darkest. The hydrogen lines
are weak, and
there are no bands from molecules. \Cref{ch1:tab:lablines} gives some laboratory lines.
(a)~Which elements are in the star's outer gas? (b)~No helium lines appear. Does this mean
the star has no helium? (c)~Suggest the star's spectral class. (d)~Find the energy gap, in
eV, that gives the \qty{589.0}{nm} line.
\begin{table}
  \caption{Laboratory wavelengths of some spectral lines.}\label{ch1:tab:lablines}
  \centering\small
  \begin{tabular}{@{}L{34mm}L{56mm}@{}}
    \toprule
    \colhead{Element} & \colhead{Lines (nm)}\\
    \midrule
    hydrogen & 410.2, 434.0, 486.1, 656.3\\
    helium & 447.1, 587.6, 667.8\\
    sodium & 589.0, 589.6\\
    ionised calcium & 393.4, 396.8\\
    \bottomrule
  \end{tabular}
\end{table}
\given{the star's dark lines; \cref{ch1:tab:lablines,ch1:tab:classes};
  $h = \qty{6.626e-34}{J.s}$, \ $c = \qty{3.00e8}{m.s^{-1}}$, \
  $\qty{1}{eV} = \qty{1.602e-19}{J}$}
\find{(a) the elements; (b) a reason; (c) the class; (d) the energy gap $E$ in eV}
\formula{(a) an element is there if its lines match the star's dark lines.\newline
  (d) $f = \dfrac{c}{\lambda}$ (from $c = f\lambda$), then $E = hf$}
\substitute{(a) Match each dark line with \cref{ch1:tab:lablines}: 393.4 and 396.8 are
  ionised calcium; 486.1 and 656.3 are hydrogen; 589.0 and 589.6 are sodium. Every dark
  line is matched.\par
  (b) Helium's visible lines are absorbed only by electrons already lifted to high
  levels. That needs a very hot star (classes O and B in \cref{ch1:tab:classes}).\par
  (c) Darkest lines from ionised calcium, metal lines (sodium) and weak hydrogen: this
  fits class G in \cref{ch1:tab:classes}.\par
  (d) Change nm to m, then use the two formulas:\par
  \eqlines{\lambda &= \qty{589.0}{nm} = \qty{5.890e-7}{m}\\
  f &= \frac{\qty{3.00e8}{m.s^{-1}}}{\qty{5.890e-7}{m}} = \qty{5.093e14}{Hz}\\
  E &= \qty{6.626e-34}{J.s} \times \qty{5.093e14}{Hz} = \qty{3.375e-19}{J}\\
  E &= \frac{\qty{3.375e-19}{J}}{\qty{1.602e-19}{J.eV^{-1}}} = \qty{2.107}{eV}
  \quad\text{(J to eV)}}}
\answer{(a) \final{hydrogen, sodium and calcium}. \ (b) No. The star is too cool to lift
  helium's electrons that high, so helium cannot absorb visible light. Stars like the Sun
  are still about one quarter helium by mass. \ (c) \final{class G}, like the Sun. \
  (d) \final{$E = \qty{2.11}{eV}$}. Sodium atoms in the cooler outer gas absorb photons
  with exactly this energy. So the star's light has a dark line at \qty{589.0}{nm}.}
\checkstep{a photon of yellow light carries about \qty{2}{eV} \cmark. No dark line is
  left without an element, so the identification is complete.}
\end{workedexample}

\begin{workedexample}{The CMB then and now}\label{ch1:we:cmb}
The CMB today matches a black body at \qty{2.725}{K}. When it was set free, the Universe
was at about \qty{3000}{K}. Find the peak wavelength then and now, and the factor by which
the waves have been stretched. How does this support an expanding Universe?
\given{$T_\mathrm{then} = \qty{3000}{K}$, \ $T_\mathrm{now} = \qty{2.725}{K}$, \
  $b = \qty{2.898e-3}{m.K}$}
\find{both values of $\lambda_\mathrm{max}$, and their ratio}
\formula{$\lambda_\mathrm{max} = \dfrac{b}{T}$}
\substitute{\eqlines{\lambda_\mathrm{then} &= \frac{\qty{2.898e-3}{m.K}}{\qty{3000}{K}} = \qty{9.66e-7}{m}\\
  \lambda_\mathrm{now} &= \frac{\qty{2.898e-3}{m.K}}{\qty{2.725}{K}} = \qty{1.063e-3}{m}}\par
  \eqlines{\text{factor} &= \frac{\qty{1.063e-3}{m}}{\qty{9.66e-7}{m}} = \num{1.10e3}}}
\answer{\final{\qty{966}{nm} (infrared) then, \qty{1.06}{mm} (microwave) now}, stretched
  \final{about 1100 times}. This supports expansion. Stretching every wavelength by the
  same factor keeps the black-body shape but lowers the temperature by that factor
  (\qty{3000}{K} $\to$ \qty{2.7}{K}). The glow arrives equally from every direction, so
  it is not one source moving away. Space itself has stretched about 1100 times while
  the light travelled.}
\checkstep{microwaves are about \qty{1}{mm} long \cmark. Also $\frac{3000}{2.725} = 1101$,
  because $\lambda_\mathrm{max} \propto \frac{1}{T}$.}
\end{workedexample}

\begin{mistakes}
\begin{itemize}
  \item \xmark\ Putting \si{\degreeCelsius} into Wien's law. \cmark\ Use kelvin.
  \item \xmark\ ``The hottest stars have the strongest hydrogen lines.'' \cmark\ Class A
        does; in O stars hydrogen is ionised.
  \item \xmark\ ``Stars of different classes are made of different elements.''
        \cmark\ Temperature changes the lines, not the mix of elements.
\end{itemize}
\end{mistakes}

\begin{examtip}
``Account for the production of emission and absorption spectra and compare them with a
black body spectrum'' (4--5 marks). Full marks need:
\begin{enumerate}
  \item Energy levels, and $E = hf$ for the one photon given out or taken in.
  \item How each kind is produced: excitation by heat or an electric current, then
        emission; or absorption by cooler gas, with the energy given out again in all
        directions.
  \item The same wavelengths for both kinds of line.
  \item \emph{Whereas} a black body gives every wavelength. Its shape and peak depend
        only on temperature (Wien's law).
\end{enumerate}
A small table is a good way to compare.
\end{examtip}

\begin{quickcheck}
  \item A sodium street lamp gives a bright yellow line at \qty{589}{nm}. Sunlight has a
        dark line at the same wavelength. Explain how each line is produced, and why both
        are at \qty{589}{nm}.
  \item A star's surface is at \qty{25000}{K}. Find $\lambda_\mathrm{max}$. In which part
        of the spectrum is it, and what is the star's class?
  \item Star P is class O9 and star Q is class M3. Both show only weak hydrogen lines.
        Give a different reason for each star.
  \item Two stars have absorption lines at the same wavelengths. Star X's lines are much
        broader than star Y's. Which star is more likely to be a giant? Explain.
\end{quickcheck}

%% ==================================================================
\section{The Hertzsprung--Russell diagram and the life of stars}\label{ch1:sec:hr}
\covers{M8.IQ1.06, M8.IQ1.07}

\begin{bigidea}
A Hertzsprung--Russell (HR) diagram plots stars by luminosity against surface
temperature. A star's place on it shows its temperature, colour, luminosity, size and
stage of life. A star's mass decides how it lives and how it dies.
\end{bigidea}

\begin{keywords}
  \kw{Luminosity ($L$)}{the total power a star gives out (W); distance does not change it.}
  \kw{Absolute magnitude ($M$)}{how bright a star would look from 10 parsecs away.
    Smaller (more negative) means brighter.}
  \kw{Main sequence}{the long stage when a star's core fuses (joins) hydrogen into helium.}
  \kw{Solar mass ($M_\odot$)}{the mass of the Sun, used as a unit for star masses.}
\end{keywords}

\subsection{Reading the HR diagram}
A street light looks bright up close and faint far away, but its power does not change.
Stars are the same.
\begin{itemize}
  \item \term{Apparent brightness} is how bright a star looks from Earth. It depends on
        the star's distance.
  \item \term{Luminosity} ($L$) is the star's true power output, in watts. We often
        compare it with the Sun's: $L_\odot = \qty{3.83e26}{W}$.
\end{itemize}

Astronomers also use the \term{magnitude} scale, a number for brightness. Smaller
numbers mean brighter. A difference of 5 magnitudes is a factor of 100.

Absolute magnitude is the magnitude a star would have from 10 parsecs away. (A
\term{parsec} (pc) is a distance unit: $\qty{1}{pc} = 3.26$ light-years.) Every star is
put at the same distance, so absolute magnitude measures luminosity. Do not mix up $M$
(absolute magnitude) with $M_\odot$ (the Sun's mass) or with class M.

\begin{figure}
  \centering
  \includegraphics{ch1_hr.pdf}
  \caption{An HR diagram (simplified). Each dashed line shows where stars with the same
    radius sit ($R_\odot$ = the Sun's radius). The numbered badges show the Sun's future:
    1 = the Sun today (main sequence), 2 = red giant, 3 = white dwarf.}\label{ch1:fig:hr}
\end{figure}

\begin{toolkit}{Reading a log scale}
On a normal scale, equal gaps mean equal \emph{added} amounts (0, 10, 20, 30). On a log
scale, equal gaps mean equal \emph{multiplying} factors (1, 10, 100, 1000). On the left
axis of \cref{ch1:fig:hr}, each labelled tick is 100 times the one below it. A star two
labelled ticks higher is $100 \times 100 = \num{10000}$ times more luminous. A log scale
lets one graph show stars from $10^{-4}$ to $10^{6}\,L_\odot$.
\end{toolkit}

Ejnar Hertzsprung and Henry Norris Russell first drew this diagram (about 1911--1913).
In \cref{ch1:fig:hr}:
\begin{itemize}
  \item \emph{Across}: surface temperature, hot on the \textbf{left}. So it also shows
        colour (blue to red) and spectral class (O to M). This axis is a log scale that
        runs backwards.
  \item \emph{Up}: luminosity (left axis) or absolute magnitude (right axis), brightest
        at the \textbf{top}. The luminosity axis is a log scale. On the magnitude axis
        the numbers get smaller towards the top, in equal steps of 5. Each step of 5 is
        a factor of 100 in luminosity.
  \item \emph{Main sequence}: the diagonal band. About \qty{90}{\percent} of stars are
        here, because it is the longest stage of a star's life.
  \item \emph{Top right}: \term{red giants} and supergiants: cool but very luminous, so
        they must be huge. \emph{Bottom left}: \term{white dwarfs}: hot but dim, so tiny.
\end{itemize}
Why size? Two stars at the same temperature give out the same power from each square
metre. So the more luminous star must have more square metres: it is bigger
(extension: $L \propto R^2T^4$).

Careful: the diagram shows \emph{surface} temperature. A red giant's core is much hotter
than the Sun's. Also, the stars of one star cluster are all at about the same distance
from us. So an HR diagram of one cluster can use apparent magnitude instead.

\begin{formulacard}[notonsheet]{Comparing luminosities from magnitudes}
  \cardformula{\frac{L_A}{L_B} = 100^{\frac{M_B - M_A}{5}}}
  \begin{symbols}
    \sym{$L_A$, $L_B$}{luminosities of stars A and B}{\si{W}}
    \sym{$M_A$, $M_B$}{absolute magnitudes of stars A and B}{no unit}
  \end{symbols}
  \inwords{every 5 magnitudes smaller means 100 times more luminous.}
  \whentouse{comparing the luminosities of two stars from their absolute magnitudes.
    (With apparent magnitudes, the same rule compares how bright they look.)}
  \whennottouse{to mix an apparent magnitude with an absolute magnitude, or to find a
    luminosity in watts.}
  \rearranged{$\dfrac{L_B}{L_A} = 100^{\frac{M_A - M_B}{5}}$ (the same rule the other way
    round). \ Useful values: 1 magnitude $= 100^{1/5} = 2.512$ times; 5 magnitudes = 100
    times; 10 magnitudes $= 10^4$ times.}
  \tensecond{$M_A = -1$, $M_B = +4$: $M_B - M_A = 5$, so A is $100^{1} = 100$ times more
    luminous.}
  \tcbbreak % layout QA (book build): page break here, so "(continued)" is not a lone Trap line
  \derive{the scale is defined so that 5 magnitudes is a factor of 100. A difference of
    $M_B - M_A$ magnitudes is $\frac{M_B - M_A}{5}$ steps of 5. Each step multiplies by
    100, so the factor is $100^{\frac{M_B - M_A}{5}}$. For example, 10 magnitudes is 2
    steps: $100 \times 100 = 100^2$.}
  \trap{a \emph{smaller} magnitude is brighter. Ten magnitudes is
    $100 \times 100 = 10^4$ times, not $100 + 100 = 200$ times.}
\end{formulacard}

\begin{workedexample}{Comparing two stars from their magnitudes}\label{ch1:we:magnitude}
Star C has absolute magnitude $+1.0$ and star D has $+8.5$. Which star is more luminous,
and by what factor?
\given{$M_C = +1.0$, \ $M_D = +8.5$}
\find{the more luminous star, and $\dfrac{L_C}{L_D}$}
\formula{$\dfrac{L_C}{L_D} = 100^{\frac{M_D - M_C}{5}}$}
\substitute{\eqlines{M_D - M_C &= 8.5 - 1.0 = 7.5\\
  &\quad\text{(C has the smaller magnitude, so C is brighter)}\\
  \frac{L_C}{L_D} &= 100^{\frac{7.5}{5}} = 100^{1.5} \quad\text{(7.5 magnitudes is 1.5
  steps of 5)}\\
  &= 100 \times 100^{0.5} = 100 \times 10 = 1000\\
  &\quad\text{(or type \texttt{100 \^{} 1.5} on a calculator)}}}
\answer{star C is \final{about $\num{1.0e3}$ times} more luminous than star D.}
\checkstep{7.5 magnitudes is between 5 ($\times 100$) and 10 ($\times\num{10000}$), and
  1000 is between these \cmark.}
\end{workedexample}

\subsection{The life of a star}
\begin{sidebyside}[0.36]
  \centering
  \begin{tikzpicture}[x=1cm, y=1cm]
    \fill[hscacc!18] (0,0) circle (1.45);
    \draw[hscacc!70, line width=0.8pt] (0,0) circle (1.45);
    \fill[hscacc!55] (0,0) circle (0.5);
    \node[font=\scriptsize, align=center, text=hscink] at (0,0) {core:\\fusion};
    \foreach \a in {45,135,225,315}
      \draw[force] (\a:2.25) -- (\a:1.55);
    \foreach \a in {0,90,180,270}
      \draw[velocity] (\a:0.6) -- (\a:1.3);
    \node[diagramlabel, text=hscforce, anchor=south] at (0,1.6) {gravity (in)};
    \node[diagramlabel, text=hscvel, align=left, anchor=west] at (1.5,-0.05) {gas\\pressure\\(out)};
  \end{tikzpicture}
  \captionof{figure}{Stellar equilibrium.}\label{ch1:fig:equilibrium}
\nextside
  \begin{itemize}
    \item Gravity pulls a cloud of gas and dust (a \term{nebula}) into clumps. A
          shrinking clump heats up: a \term{protostar}.
    \item At about 10 million K, \term{nuclear fusion} starts in the core. Hydrogen
          nuclei join to make helium, and this gives out energy
          (\cref{ch1:sec:fusion}). The star is now on the main sequence.
    \item Gravity pulls in; hot gas (and radiation) pushes out. They balance:
          \term{stellar equilibrium} (\cref{ch1:fig:equilibrium}).
  \end{itemize}
\end{sidebyside}

A more massive star has stronger gravity, a hotter core and much faster fusion. It is
hotter, bluer and more luminous, but its life is much shorter. The Sun lasts about 10
billion years, and a $10\,M_\odot$ star only about 30 million years. A $0.1\,M_\odot$
star lasts trillions of years.

When the core runs out of hydrogen, it shrinks and heats up, and hydrogen fuses in a
shell around it. The outer layers swell and cool: a red giant (point 2 in
\cref{ch1:fig:hr}). At about 100 million K, helium fuses into carbon.
\Cref{ch1:fig:lifecycle} shows the two possible endings.

\begin{figure}
  \centering
  \begin{tikzpicture}[x=1mm, y=1mm,
      st/.style={concept=#1, font=\footnotesize, inner sep=1.4mm},
      fin/.style={concept=#1, font=\footnotesize, inner sep=1.4mm, fill=hscbrand!16}]
    \node[st=12mm] (neb) at (7.4,0) {nebula};
    \node[st=15mm] (pro) at (27.9,0) {protostar};
    \node[st=15mm] (ms) at (49.9,0) {main\\sequence};
    \node[st=17mm] (rg) at (76.9,11) {red giant};
    \node[st=17mm] (rsg) at (76.9,-11) {red\\supergiant};
    \node[st=15mm] (pn) at (99.9,11) {planetary\\nebula};
    \node[st=15mm] (sn) at (99.9,-11) {supernova};
    \node[fin=24mm] (wd) at (126.4,11) {white dwarf\\(core below $1.4\,M_\odot$)};
    \node[fin=24mm] (ns) at (126.4,-4) {neutron star\\(core $1.4$--$3\,M_\odot$)};
    \node[fin=24mm] (bh) at (126.4,-18) {black hole\\(core above $3\,M_\odot$)};
    \draw[link] (neb) -- (pro);
    \draw[link] (pro) -- (ms);
    \draw[link] (ms.east) -- (rg.west);
    \draw[link] (ms.east) -- (rsg.west);
    \node[font=\footnotesize\bfseries, text=hscbrand, anchor=south west] at (64,17.5)
      {Starting mass below $8\,M_\odot$};
    \node[font=\footnotesize\bfseries, text=hscbrand, anchor=north west] at (64,-17.5)
      {Starting mass above $8\,M_\odot$};
    \draw[link] (rg) -- (pn);
    \draw[link] (pn) -- (wd);
    \draw[link] (rsg) -- (sn);
    \draw[link] (sn.east) -- (ns.west);
    \draw[link] (sn.east) -- (bh.west);
  \end{tikzpicture}
  \caption{The life cycle of a star depends on its starting mass. The masses on the right
    are those of the leftover core.}\label{ch1:fig:lifecycle}
\end{figure}

\begin{itemize}
  \item \emph{Below about $8\,M_\odot$} (like the Sun): the outer layers drift away as a
        \term{planetary nebula}. The hot core is left as a white dwarf, about the size of
        Earth. With no fusion it slowly cools. (A cold ``black dwarf'' would take longer
        than the age of the Universe, so none exists yet.)
  \item \emph{Above about $8\,M_\odot$}: the star swells into a \term{red supergiant}, a
        huge, cool, very luminous star. Its core fuses heavier elements, up to iron.
        Then the iron core collapses and the star explodes. This explosion is a
        \term{supernova}. The leftover core becomes a \term{neutron star}: a ball about
        \qty{20}{km} across, made almost only of neutrons. Its mass is 1.4 to 3 times the
        Sun's mass. If the core is heavier than about $3\,M_\odot$, it becomes a
        \term{black hole}: its gravity is so strong that nothing, not even light, can
        escape.
\end{itemize}

\begin{workedexample}{Reading the HR diagram}\label{ch1:we:hr}
Star P: $T = \qty{3300}{K}$, $L = 10^5\,L_\odot$. Star Q: $T = \qty{20000}{K}$,
$L = 10^{-2}\,L_\odot$. Use \cref{ch1:fig:hr} to give each star's colour, class, type
and size. Find each star's peak wavelength.
\given{P: \qty{3300}{K}, $10^5\,L_\odot$; \ Q: \qty{20000}{K}, $10^{-2}\,L_\odot$; \
  $b = \qty{2.898e-3}{m.K}$}
\find{colour, class, type and size from the diagram; $\lambda_\mathrm{max}$}
\formula{read the diagram for colour, class, type and size; then
  $\lambda_\mathrm{max} = \dfrac{b}{T}$}
\substitute{P: find \qty{3300}{K} on the bottom axis. It is in the M band, so P is red,
  class M. Go up to $10^5\,L_\odot$ on the left axis. That point is in the Supergiants
  band, at the top right. The nearest dashed line is $1000\,R_\odot$.\newline
  Q: find \qty{20000}{K}. It is in the B band, so Q is blue-white, class B. Go up to
  $10^{-2}\,L_\odot$. That point is in the White dwarfs region, at the bottom left, near
  the $0.01\,R_\odot$ line.\newline
  \eqlines{\text{P: } \lambda_\mathrm{max} &= \frac{\qty{2.898e-3}{m.K}}{\qty{3300}{K}}
  = \qty{8.78e-7}{m} = \qty{878}{nm}\\
  \text{Q: } \lambda_\mathrm{max} &= \frac{\qty{2.898e-3}{m.K}}{\qty{20000}{K}}
  = \qty{1.45e-7}{m} = \qty{145}{nm}}}
\answer{P: red, class M, top right: a \final{red supergiant}, about $1000\,R_\odot$; peak
  \qty{878}{nm} (infrared). Q: blue-white, class B, bottom left: a \final{white dwarf},
  about $0.01\,R_\odot$; peak \qty{145}{nm} (ultraviolet).}
\checkstep{P is cool but very luminous, so huge; Q is hot but dim, so tiny \cmark.}
\end{workedexample}

\begin{workedexample}{Two stars, two endings}\label{ch1:we:endings}
Star R starts with $2\,M_\odot$ and star S starts with $15\,M_\odot$. Describe each
star's life and name its final remnant.
\given{starting masses $2\,M_\odot$ (R) and $15\,M_\odot$ (S)}
\find{the stages and the remnant of each star}
\formula{the rule in \cref{ch1:fig:lifecycle}: starting mass below about $8\,M_\odot$
  $\to$ red giant $\to$ planetary nebula $\to$ white dwarf; above about $8\,M_\odot$
  $\to$ red supergiant $\to$ supernova $\to$ neutron star or black hole}
\substitute{R: $2 < 8$, so R follows the top path. \ S: $15 > 8$, so S follows the
  bottom path.}
\answer{Both: nebula $\to$ protostar $\to$ main sequence (H to He).\newline
  R: about 1--2 billion years on the main sequence $\to$ red giant (He to C)
  $\to$ planetary nebula $\to$ \final{white dwarf}.\newline
  S: hotter and far more luminous, only about 10--15 million years on the main
  sequence $\to$ red supergiant (fusion up to iron) $\to$ supernova $\to$
  \final{neutron star or black hole} (this depends on the mass of the leftover core).}
\checkstep{S has more fuel, but burns it much faster, so it dies sooner \cmark.}
\end{workedexample}

\begin{mistakes}
\begin{itemize}
  \item \xmark\ Temperature increasing to the right. \cmark\ Hot stars are on the left.
  \item \xmark\ ``A bigger magnitude is brighter.'' \cmark\ Smaller (more negative) is brighter.
  \item \xmark\ ``A red supergiant becomes a white dwarf.'' \cmark\ It explodes as a
        supernova: neutron star or black hole.
\end{itemize}
\end{mistakes}

\begin{examtip}
``Explain how an HR diagram shows a star's surface temperature, colour and luminosity''
(5--6 marks). Full marks need:
\begin{enumerate}
  \item The horizontal axis: surface temperature, hot on the left, log scale.
  \item Colour comes from temperature: hot = blue, cool = red (Wien's law).
  \item The vertical axis: luminosity or absolute magnitude, brightest at the top.
  \item What each group means and its stage of life: main sequence; giants (cool but
        bright, so big); white dwarfs (hot but dim, so small).
  \item A neat labelled sketch.
\end{enumerate}
``Path of a Sun-like star'': main sequence $\to$ up and right (red giant) $\to$ left and
down (white dwarf).
\end{examtip}

\begin{quickcheck}
  \item Where on an HR diagram are stars that are cool but very luminous? What does this
        tell you about their size?
  \item Star A has absolute magnitude $-5$ and star B $+10$. Which is more luminous, and
        by about what factor?
  \item Explain why a $10\,M_\odot$ star spends much less time on the main sequence than
        the Sun.
\end{quickcheck}

%% ==================================================================
\section{Fusion in stars: \texorpdfstring{$E = mc^2$}{E = mc2} and nucleosynthesis}\label{ch1:sec:fusion}
\covers{M8.IQ1.03, M8.IQ1.07}

\begin{bigidea}
Stars shine by fusing light nuclei into heavier ones. The products have a little less
mass, and the lost mass becomes energy: $E = \Delta m\,c^2$. Main sequence stars fuse
hydrogen into helium. Stars that have left the main sequence (giants and supergiants)
build heavier elements, up to iron.
\end{bigidea}

\begin{keywords}
  \kw{Nuclear fusion}{two light nuclei joining to make one heavier nucleus.}
  \kw{Mass defect ($\Delta m$)}{the mass that disappears in a nuclear reaction: mass before
    minus mass after.}
  \kw{Catalyst}{something that helps a reaction but is not used up.}
\end{keywords}

\begin{refresher}{Nuclear symbols}
$\chOneNuc{A}{Z}{X}$: $Z$ = atomic number (protons), $A$ = mass number (protons +
neutrons). Proton $\chOneNuc{1}{1}{H}$, positron $\chOneNuc{0}{+1}{e}$, gamma photon
$\gamma$. \emph{Deuterium} $\chOneNuc{2}{1}{H}$ is hydrogen with one proton and one
neutron. A \emph{neutrino} ($\nu_\mathrm{e}$) is a tiny neutral particle with almost no
mass. In every nuclear equation, the $A$ values on the left add up to the same total as
on the right. The $Z$ values do too.
\end{refresher}

\subsection{Mass becomes energy}
Four hydrogen nuclei have a little more mass than one helium nucleus. When they fuse,
about \qty{0.7}{\percent} of the mass disappears and comes out as energy.

Fusion needs more than about 10 million K. Nuclei are positive, so they repel each
other. Only very fast (very hot) nuclei get close enough to join. Fusion gives out
energy only for nuclei up to iron (\cref{ch4:sec:binding} shows why).

Fusing \qty{1}{kg} of hydrogen turns \qty{0.00712}{kg} into energy:
$E = mc^2 = \qty{0.00712}{kg} \times (\qty{3.00e8}{m.s^{-1}})^2 = \qty{6.4e14}{J}$.
Burning \qty{1}{kg} of petrol gives only about \qty{4.6e7}{J}. So fusion gives more than
10 million times as much ($\num{6.4e14} \div \num{4.6e7} = \num{1.4e7}$). That is why
stars can shine for billions of years.

Nuclei are so light that we measure them in \term{atomic mass units} (u):
$\qty{1}{u} = \qty{1.661e-27}{kg}$, about the mass of one proton. The data sheet also
gives $\qty{1}{u} = \qty{931.5}{MeV}/c^2$. This means \qty{1}{u} of mass, turned fully
into energy, gives \qty{931.5}{MeV}.

\begin{formulacard}[combine={$E = mc^2$ and $\mathrm{u} = 931.5\ \mathrm{MeV}/c^2$}]{Energy released in a nuclear reaction}
  \cardformula{E = \Delta m\,c^2 \qquad E\,(\mathrm{MeV}) = \Delta m\,(\mathrm{u}) \times 931.5}
  \begin{symbols}
    \sym{$E$}{energy released}{\si{J} (or \si{MeV})}
    \sym{$\Delta m$}{mass defect = mass before $-$ mass after; in u, where
      $\qty{1}{u} = \qty{1.661e-27}{kg} = \qty{931.5}{MeV}/c^2$}{\si{kg} (or u)}
    \sym{$c$}{speed of light, \qty{3.00e8}{m.s^{-1}}}{\si{m.s^{-1}}}
    \sym{$P$, $\Delta t$}{power given out (a star's luminosity); time taken}{\si{W}; \si{s}}
  \end{symbols}
  \inwords{the mass that disappears in a reaction comes out as energy.}
  \whentouse{fusion and other nuclear reactions with known masses; a star's mass loss.}
  \whennottouse{with mass numbers ($A$): they always balance, giving $\Delta m = 0$.}
  \rearranged{$\Delta m = \dfrac{E}{c^2}$ \qquad mass lost each second: power $P$ is
    energy per second, so divide $E = \Delta m\,c^2$ by the time:
    $\dfrac{\Delta m}{\Delta t} = \dfrac{P}{c^2}$}
  \tensecond{$\Delta m = \qty{0.0200}{u}$: $E = 0.0200 \times 931.5 = \qty{18.6}{MeV}$.}
  \derive{with $\qty{1}{u} = \qty{931.5}{MeV}/c^2$ (data sheet):
    $E = \Delta m\,c^2 = \Delta m\,(\mathrm{u}) \times \dfrac{\qty{931.5}{MeV}}{c^2} \times c^2
    = \Delta m\,(\mathrm{u}) \times \qty{931.5}{MeV}$.}
  \trap{use one route only: u $\to$ $\times\,931.5$ MeV, or kg $\to$ $\times\,c^2$ for J.
    Keep every decimal place of the masses.}
\end{formulacard}

\subsection{Main sequence stars: the proton--proton chain}
Stars like the Sun (up to about $1.3\,M_\odot$, core below about 18 million K) mainly use
the \term{proton--proton chain}. It is a chain of three fusion steps that turns four
protons into one helium-4 nucleus (\cref{ch1:fig:pp}). Helium-4 has four particles (2
protons + 2 neutrons). But four protons almost never hit each other at the same moment.
So the star builds helium two particles at a time.

\begin{figure}
  \centering
  \begin{tikzpicture}[x=1mm, y=1mm,
      rx/.style={concept=54mm, font=\small, inner sep=1.5mm},
      st/.style={font=\footnotesize\bfseries, text=hscbrand, anchor=east}]
    \node[rx] (a1) at (-31,0)
      {$\chOneNuc{1}{1}{H} + \chOneNuc{1}{1}{H} \to \chOneNuc{2}{1}{H} + \chOneNuc{0}{+1}{e} + \nu_\mathrm{e}$};
    \node[rx] (a2) at (31,0)
      {$\chOneNuc{1}{1}{H} + \chOneNuc{1}{1}{H} \to \chOneNuc{2}{1}{H} + \chOneNuc{0}{+1}{e} + \nu_\mathrm{e}$};
    \node[rx] (b1) at (-31,-13) {$\chOneNuc{2}{1}{H} + \chOneNuc{1}{1}{H} \to \chOneNuc{3}{2}{He} + \gamma$};
    \node[rx] (b2) at (31,-13) {$\chOneNuc{2}{1}{H} + \chOneNuc{1}{1}{H} \to \chOneNuc{3}{2}{He} + \gamma$};
    \node[rx, fill=hscbrand!16] (c) at (0,-27)
      {$\chOneNuc{3}{2}{He} + \chOneNuc{3}{2}{He} \to \chOneNuc{4}{2}{He} + \chOneNuc{1}{1}{H} + \chOneNuc{1}{1}{H}$};
    \draw[link] (a1) -- (b1);
    \draw[link] (a2) -- (b2);
    \draw[link] (b1.south) -- (c.north west);
    \draw[link] (b2.south) -- (c.north east);
    \node[st] at (-61,0) {Step 1};
    \node[st] at (-61,-13) {Step 2};
    \node[st] at (-61,-27) {Step 3};
  \end{tikzpicture}
  \caption{The proton--proton chain. In step 1 a proton turns into a neutron. Steps 1 and
    2 happen twice for every step 3, which gives two protons back.}\label{ch1:fig:pp}
\end{figure}

Count what goes in and out to make one helium-4 nucleus:
\begin{itemize}
  \item \emph{In}: steps 1 and 2 each happen twice. That is $4\,\chOneNuc{1}{1}{H}$
        (step 1) $+\ 2\,\chOneNuc{1}{1}{H}$ (step 2) $= 6\,\chOneNuc{1}{1}{H}$.
  \item \emph{Made and then used again} (these cancel): $2\,\chOneNuc{2}{1}{H}$ (made in
        step 1, used in step 2) and $2\,\chOneNuc{3}{2}{He}$ (made in step 2, used in
        step 3).
  \item \emph{Out}: $\chOneNuc{4}{2}{He} + 2\,\chOneNuc{1}{1}{H}$ (step 3),
        $2\,\chOneNuc{0}{+1}{e} + 2\nu_\mathrm{e}$ (step 1, twice), $2\gamma$ (step 2,
        twice).
  \item 6 protons in and 2 protons out means 4 protons are used.
\end{itemize}
So the net change is
\[
  4\,\chOneNuc{1}{1}{H} \to \chOneNuc{4}{2}{He} + 2\,\chOneNuc{0}{+1}{e} + 2\nu_\mathrm{e} + 2\gamma
  \qquad (\qty{24.7}{MeV}).
\]
The two positrons soon annihilate with two electrons. This gives more gamma photons and
about \qty{2.0}{MeV} more energy. So the total is about \qty{26.7}{MeV} for each helium-4
nucleus.

\subsection{Main sequence stars: the CNO cycle}
Stars above about $1.3\,M_\odot$ have hotter cores (above about 18 million K). They
mainly use the \term{CNO cycle} (\cref{ch1:fig:cno}): a carbon-12 nucleus catches protons
one at a time.

\begin{sidebyside}[0.46]
  \centering
  \begin{tikzpicture}[x=1cm, y=1cm,
      nuc/.style={circle, draw=hscbrand!60, fill=hscbrand!9, line width=0.7pt,
        minimum size=8.5mm, inner sep=0pt, font=\small},
      lab/.style={font=\scriptsize, text=hscmuted, align=center}]
    \def\R{1.85}
    \foreach \ang/\name/\txt in {90/c12/{\chOneNuc{12}{6}{C}}, 30/n13/{\chOneNuc{13}{7}{N}},
        -30/c13/{\chOneNuc{13}{6}{C}}, -90/n14/{\chOneNuc{14}{7}{N}},
        -150/o15/{\chOneNuc{15}{8}{O}}, 150/n15/{\chOneNuc{15}{7}{N}}}
      \node[nuc] (\name) at (\ang:\R) {$\txt$};
    \foreach \a/\b/\pos/\txt in {c12/n13/60/{1: $+\,^{1}\mathrm{H}$\\$\gamma$ out},
        n13/c13/0/{2: $\mathrm{e}^{+}$, $\nu$\\out},
        c13/n14/-60/{3: $+\,^{1}\mathrm{H}$\\$\gamma$ out},
        n14/o15/-120/{4: $+\,^{1}\mathrm{H}$\\$\gamma$ out},
        o15/n15/180/{5: $\mathrm{e}^{+}$, $\nu$\\out},
        n15/c12/120/{6: $+\,^{1}\mathrm{H}$\\$^{4}\mathrm{He}$ out}}{
      \draw[link, line width=0.8pt] (\a) to[bend left=22] (\b);
      \node[lab] at (\pos:2.72) {\txt};}
    \node[font=\footnotesize, align=center, text=hscink] at (0,0)
      {$^{12}\mathrm{C}$ is a\\catalyst:\\used and\\returned};
  \end{tikzpicture}
  \captionof{figure}{The CNO cycle turns four protons into one helium-4
    nucleus.}\label{ch1:fig:cno}
\nextside
  \small
  \begin{enumerate}
    \item $\chOneNuc{12}{6}{C} + \chOneNuc{1}{1}{H} \to \chOneNuc{13}{7}{N} + \gamma$
    \item $\chOneNuc{13}{7}{N} \to \chOneNuc{13}{6}{C} + \chOneNuc{0}{+1}{e} + \nu_\mathrm{e}$
    \item $\chOneNuc{13}{6}{C} + \chOneNuc{1}{1}{H} \to \chOneNuc{14}{7}{N} + \gamma$
    \item $\chOneNuc{14}{7}{N} + \chOneNuc{1}{1}{H} \to \chOneNuc{15}{8}{O} + \gamma$
    \item $\chOneNuc{15}{8}{O} \to \chOneNuc{15}{7}{N} + \chOneNuc{0}{+1}{e} + \nu_\mathrm{e}$
    \item $\chOneNuc{15}{7}{N} + \chOneNuc{1}{1}{H} \to \chOneNuc{12}{6}{C} + \chOneNuc{4}{2}{He}$
  \end{enumerate}
  Count: 4 protons go in (steps 1, 3, 4, 6). 2 positrons and 2 neutrinos come out
  (steps 2, 5). 3 gammas come out (steps 1, 3, 4). Carbon-12 comes back at the end.\par
  Net: $4\,\chOneNuc{1}{1}{H} \to \chOneNuc{4}{2}{He} + 2\,\chOneNuc{0}{+1}{e}
  + 2\nu_\mathrm{e} + 3\gamma$
\end{sidebyside}

The carbon-12 used in step 1 comes back in step 6, so it is a \term{catalyst}. But the
star must already contain some carbon, made by earlier stars. \Cref{ch1:tab:ppcno}
compares the two routes.

\begin{table}
  \caption{The proton--proton chain and the CNO cycle.}\label{ch1:tab:ppcno}
  \centering\small
  \begin{tabular}{@{}L{30mm}L{53mm}L{55mm}@{}}
    \toprule
    & \colhead{Proton--proton chain} & \colhead{CNO cycle}\\
    \midrule
    Main route in & the Sun and smaller stars (below about $1.3\,M_\odot$) &
      more massive stars (above about $1.3\,M_\odot$)\\
    Core temperature & below about 18 million K & above about 18 million K\\
    Needs & only hydrogen & carbon, nitrogen or oxygen already there (catalysts)\\
    Net result & 4 protons $\to$ helium-4, \qty{26.7}{MeV} & the same\\
    \bottomrule
  \end{tabular}
\end{table}

\subsection{After the main sequence: helium and heavier elements}
In a red giant the core reaches about 100 million K, and helium fuses into carbon by
the \term{triple-alpha process} (helium-4 nuclei are alpha particles):
\[
  \chOneNuc{4}{2}{He} + \chOneNuc{4}{2}{He} \to \chOneNuc{8}{4}{Be} \qquad
  \chOneNuc{8}{4}{Be} + \chOneNuc{4}{2}{He} \to \chOneNuc{12}{6}{C} + \gamma \qquad
  \text{net: } 3\,\chOneNuc{4}{2}{He} \to \chOneNuc{12}{6}{C} + \text{energy}
\]
Beryllium-8 breaks apart almost at once, so a third helium nucleus must arrive very
quickly. Some carbon then catches another helium nucleus:
$\chOneNuc{12}{6}{C} + \chOneNuc{4}{2}{He} \to \chOneNuc{16}{8}{O} + \gamma$.

\begin{sidebyside}[0.3]
  \centering
  \begin{tikzpicture}[x=1cm, y=1cm]
    \foreach \r/\c in {2.0/hscacc!12, 1.62/hscacc!22, 1.27/hscacc!32, 0.97/hscacc!42,
        0.72/hscacc!52, 0.48/hscmuted!35}
      \fill[\c] (0,0) circle (\r);
    \fill[hscmuted!70] (0,0) circle (0.25);
    \foreach \r in {2.0, 1.62, 1.27, 0.97, 0.72, 0.48, 0.25}
      \draw[hscink!40, line width=0.4pt] (0,0) circle (\r);
    \foreach \y/\t in {1.81/H, 1.445/He, 1.12/C, 0.845/Ne, 0.6/O, 0.365/Si}
      \node[font=\scriptsize, text=hscink] at (0,\y) {\t};
    \node[font=\scriptsize\bfseries, text=white] at (0,0) {Fe};
  \end{tikzpicture}
  \captionof{figure}{Fusion shells in a massive star (not to scale).}\label{ch1:fig:onion}
\nextside
  \begin{itemize}
    \item In massive stars the core keeps getting hotter. It fuses carbon, then neon,
          oxygen and silicon. Each fuel burns in its own shell, like the layers of an
          onion (\cref{ch1:fig:onion}).
    \item Fusing iron \emph{takes in} energy. With no energy source left, the iron core
          collapses: a supernova.
    \item Elements heavier than iron (gold, uranium) form when nuclei catch neutrons:
          in supernovae, neutron-star collisions and old giant stars.
    \item Planetary nebulae and supernovae spread these elements into space, where new
          stars and planets form.
  \end{itemize}
\end{sidebyside}

So the elements came from three places. The Big Bang made hydrogen, most of the helium
and a little lithium. Fusion in stars made more helium, and carbon up to iron. Neutron
capture made the elements heavier than iron. The hydrogen in your body (in water, for
example) was made in the Big Bang. Almost every heavier atom in your body, such as
carbon, oxygen, calcium and iron, was made in stars.

\begin{workedexample}{Energy from hydrogen fusion}\label{ch1:we:ppenergy}
Use the atomic masses $m(\chOneNuc{1}{1}{H}) = \qty{1.007825}{u}$ and
$m(\chOneNuc{4}{2}{He}) = \qty{4.002603}{u}$. Find the total energy released when four
hydrogen nuclei fuse into one helium nucleus (the whole proton--proton chain). Also find
the percentage of the mass that becomes energy.
\given{$m(\mathrm{H}) = \qty{1.007825}{u}$, \ $m(\mathrm{He}) = \qty{4.002603}{u}$, \
  $\qty{1}{u} = \qty{931.5}{MeV}/c^2$, \ $\qty{1}{eV} = \qty{1.602e-19}{J}$}
\find{$E$ in MeV and in J, and $\dfrac{\Delta m}{m_\mathrm{before}} \times \qty{100}{\percent}$}
\formula{$\Delta m = 4m(\mathrm{H}) - m(\mathrm{He})$, \ then
  $E\,(\mathrm{MeV}) = \Delta m\,(\mathrm{u}) \times 931.5$}
\substitute{\eqlines{\Delta m &= 4 \times \qty{1.007825}{u} - \qty{4.002603}{u}\\
  &= \qty{4.031300}{u} - \qty{4.002603}{u} = \qty{0.028697}{u}\\
  E &= \num{0.028697} \times \qty{931.5}{MeV} = \qty{26.73}{MeV}}\par
  \eqlines{E &= \qty{26.73e6}{eV} \times \qty{1.602e-19}{J.eV^{-1}} \quad\text{(MeV to J)}\\
  &= \qty{4.282e-12}{J}\\
  \text{percentage} &= \frac{\qty{0.028697}{u}}{\qty{4.031300}{u}} \times \qty{100}{\percent}
  = \qty{0.712}{\percent}}}
\answer{\final{$E = \qty{26.7}{MeV} = \qty{4.28e-12}{J}$}; \ \final{\qty{0.712}{\percent}}
  of the mass}
\checkstep{matches the \qty{26.7}{MeV} and ``\qty{0.7}{\percent}'' above \cmark. These
  masses are for whole atoms, so they include the electrons. So this answer also
  includes the energy from the two positrons annihilating with electrons.}
\end{workedexample}

\begin{workedexample}{Does a star die because its mass runs out?}\label{ch1:we:massloss}
Sirius A is a main sequence star. Its luminosity is \qty{9.73e27}{W} and its mass is
\qty{4.10e30}{kg}. (a) How much mass does it turn into energy each second? (b) What
percentage of its mass is that in 1.0 billion years? (c) A student says: ``A star dies
because fusion has turned most of its mass into energy.'' Use (b) to evaluate this.
\given{$P = \qty{9.73e27}{W}$, \ $m_\mathrm{star} = \qty{4.10e30}{kg}$, \
  $c = \qty{3.00e8}{m.s^{-1}}$,\newline
  $t = \num{1.0e9}\ \text{years} \times \qty{3.156e7}{s.year^{-1}} = \qty{3.156e16}{s}$}
\find{(a) $\dfrac{\Delta m}{\Delta t}$; \ (b) the percentage of $m_\mathrm{star}$; \
  (c) a judgement, using (b)}
\formula{$\dfrac{\Delta m}{\Delta t} = \dfrac{P}{c^2}$, \ then mass $=$ rate $\times$ time}
\substitute{\eqlines{\frac{\Delta m}{\Delta t} &= \frac{\qty{9.73e27}{W}}{(\qty{3.00e8}{m.s^{-1}})^2}
  = \qty{1.081e11}{kg.s^{-1}}}\par
  \eqlines{\text{mass} &= \qty{1.081e11}{kg.s^{-1}} \times \qty{3.156e16}{s} = \qty{3.412e27}{kg}\\
  \text{percentage} &= \frac{\qty{3.412e27}{kg}}{\qty{4.10e30}{kg}} \times \qty{100}{\percent}
  = \qty{0.0832}{\percent}}}
\answer{(a) \final{\qty{1.08e11}{kg} each second} \ (b) \final{only \qty{0.0832}{\percent}}
  \newline (c) The statement is wrong. In a billion years, fusion turns less than
  \qty{0.1}{\percent} of the star's mass into energy (part b). A star leaves the main
  sequence when the hydrogen in its \emph{core} runs out, not when its mass runs out.}
\checkstep{\qty{1.08e11}{kg} each second sounds huge, but it is tiny next to the star's
  mass \cmark.}
\end{workedexample}

\begin{mistakes}
\begin{itemize}
  \item \xmark\ Using mass numbers: $4 \times 1 - 4 = 0$. \cmark\ Use the real masses, with
        every decimal place.
  \item \xmark\ ``Carbon is used up in the CNO cycle.'' \cmark\ It comes back: a catalyst.
  \item \xmark\ Leaving out the positron and neutrino in step 1. \cmark\ Check that $Z$
        balances: $1 + 1 = 1 + 1$.
\end{itemize}
\end{mistakes}

\begin{examtip}
``Compare the proton--proton chain and the CNO cycle'' (3--4 marks). Full marks need:
\begin{enumerate}
  \item One similarity: both turn 4 protons into one helium-4 nucleus and release
        \qty{26.7}{MeV}, in main sequence stars.
  \item Difference 1: p--p in the Sun and smaller stars (core below about 18 million K);
        CNO in more massive, hotter stars.
  \item Difference 2: CNO needs carbon already in the star, as a catalyst. The p--p
        chain needs only hydrogen.
\end{enumerate}
Write balanced equations with $A$ and $Z$.

For ``analyse $E = mc^2$ in stars'', link mass defect $\to$ energy $\to$ huge energy per
kilogram $\to$ billions of years of shining.
\end{examtip}

\begin{quickcheck}
  \item With $m(\chOneNuc{4}{2}{He}) = \qty{4.002603}{u}$ and $m(\chOneNuc{12}{6}{C}) =
        \qty{12.000000}{u}$, find the energy (MeV) released when three helium-4 nuclei
        fuse into carbon-12.
  \item Complete this CNO step: $\chOneNuc{13}{6}{C} + \chOneNuc{1}{1}{H} \to \ ? \ + \gamma$.
  \item Give two differences between the proton--proton chain and the CNO cycle.
  \item Explain why elements heavier than iron are not made by fusion in a star's core.
\end{quickcheck}

%% ==================================================================
\section{Key experiments}\label{ch1:sec:experiments}
\covers{M8.IQ1.02, M8.IQ1.04, M8.IQ1.05, M8.IQ1.06}

\begin{experiment}{Hubble's law from galaxy data}
\aim{To show that recession velocity is proportional to distance, and to estimate $H_0$
  and the age of the Universe.}
\apparatus{Hubble (1929): the \qty{2.5}{m} telescope at Mount Wilson with a spectrograph
  (an instrument that spreads light into a spectrum and records it). Each galaxy's
  spectrum was photographed next to a laboratory lamp's spectrum
  (\cref{ch1:fig:hubbleapp}). Speeds came from redshifts measured by Vesto Slipher and
  by Milton Humason, Hubble's assistant. Distances came from Cepheid variables. In
  class: a table of galaxy data and a spreadsheet.
  \begin{figure}
    \centering
    \begin{tikzpicture}[x=1cm, y=1cm,
        box/.style={draw=hscbrand!70, fill=hscbrand!8, rounded corners=1pt,
          minimum height=8.5mm, align=center, font=\footnotesize, inner sep=2pt}]
      \fill[hscmuted!35] (0,0) ellipse (0.55 and 0.22);
      \fill[hscmuted!75] (0,0) ellipse (0.17 and 0.08);
      \node[diagramlabel, below, font=\footnotesize] at (0,-0.3) {galaxy};
      \node[box, minimum width=21mm] (tel) at (2.55,0) {telescope\\(\qty{2.5}{m} mirror)};
      \node[box, minimum width=9mm] (slit) at (4.65,0) {slit};
      \node[box, minimum width=15mm] (gr) at (6.25,0) {prism or\\grating};
      \node[box, minimum width=20mm] (plate) at (8.6,0) {photographic\\plate};
      \node[box, minimum width=20mm] (lamp) at (4.65,1.45) {laboratory lamp};
      \draw[ray] (0.6,0) -- (tel.west);
      \draw[ray] (tel) -- (slit);
      \draw[ray] (slit) -- (gr);
      \draw[ray] (gr) -- (plate);
      \draw[ray] (lamp) -- (slit);
      \draw[hscmuted, decorate, decoration={brace, amplitude=4pt, mirror}]
        (4.15,-0.6) -- (9.65,-0.6) node[midway, below=5pt, font=\footnotesize,
        text=hscmuted] {spectrograph};
      \begin{scope}[shift={(2.2,-3.35)}]
        \fill[hscink!85] (0,1.0) rectangle (6,1.35);
        \fill[hscacc!14] (0,0.45) rectangle (6,0.9);
        \fill[hscink!85] (0,0) rectangle (6,0.35);
        \foreach \x in {1.0, 2.9, 4.6} {
          \draw[hscacc!70!yellow, line width=1.8pt] (\x,1.0) -- (\x,1.35);
          \draw[hscacc!70!yellow, line width=1.8pt] (\x,0) -- (\x,0.35);
          \draw[hscink, line width=1.5pt] (\x+0.4,0.45) -- (\x+0.4,0.9);}
        \node[font=\footnotesize, anchor=east] at (-0.1,1.175) {lamp};
        \node[font=\footnotesize, anchor=east] at (-0.1,0.675) {galaxy};
        \node[font=\footnotesize, anchor=east] at (-0.1,0.175) {lamp};
        \draw[guide] (4.6,0) -- (4.6,-0.3);
        \draw[guide] (5.0,0.45) -- (5.0,-0.3);
        \draw[hscink, {Stealth[length=1.6mm]}-{Stealth[length=1.6mm]}] (4.6,-0.25)
          -- (5.0,-0.25) node[right, font=\footnotesize] {$\Delta\lambda$};
        \node[font=\footnotesize, text=hscmuted, anchor=north] at (3,-0.35)
          {longer wavelength (redder) $\rightarrow$};
      \end{scope}
    \end{tikzpicture}
    \caption{How redshifts were measured (not to scale). The galaxy's spectrum was
      photographed between two spectra of a laboratory lamp. Each galaxy line sits a
      distance $\Delta\lambda$ to the red side of the matching lamp
      line.}\label{ch1:fig:hubbleapp}
  \end{figure}}
\method{\begin{enumerate}
  \item For each galaxy, find $\Delta\lambda = \lambda - \lambda_0$ for the H-alpha line
        ($\lambda_0 = \qty{656.3}{nm}$), then $v = \dfrac{c\,\Delta\lambda}{\lambda_0}$ in
        \si{km.s^{-1}}.
  \item Plot $v$ against $d$. Draw the line of best fit through the origin.
  \item Gradient $= H_0$; then $t \approx \dfrac{1}{H_0}$ (with $H_0$ in \si{s^{-1}}).
\end{enumerate}}
\results{{\centering\small
  \begin{tabular}{@{}lccccc@{}}
    \toprule
    Galaxy & A & B & C & D & E\\
    \midrule
    Distance $d$ (\si{Mpc}) & 22 & 48 & 75 & 104 & 140\\
    Observed $\lambda$ (\si{nm}) & 659.8 & 663.3 & 668.1 & 671.7 & 678.0\\
    $\Delta\lambda$ (\si{nm}) & 3.5 & 7.0 & 11.8 & 15.4 & 21.7\\
    $v$ (\si{km.s^{-1}}) & \num{1600} & \num{3200} & \num{5390} & \num{7040} & \num{9920}\\
    \bottomrule
  \end{tabular}\par}\medskip
  % (a group, not a center environment, so the Results heading stays with its table)
  Sample working (galaxy A):
  $v = \dfrac{c\,\Delta\lambda}{\lambda_0} = \dfrac{\qty{3.00e5}{km.s^{-1}} \times
  \qty{3.5}{nm}}{\qty{656.3}{nm}} = \qty{1.6e3}{km.s^{-1}}$ (the nm cancel).
  The graph is \cref{ch1:fig:hubble}: gradient \qty{69.9}{km.s^{-1}.Mpc^{-1}}.}
\conclusion{$v \propto d$, with $H_0 \approx \qty{70}{km.s^{-1}.Mpc^{-1}}$. The age:
  $H_0 = \dfrac{\qty{69.9}{km.s^{-1}}}{\qty{3.086e19}{km}} = \qty{2.265e-18}{s^{-1}}$, so
  $t = \dfrac{1}{H_0} = \qty{4.415e17}{s}$, and
  $\dfrac{\qty{4.415e17}{s}}{\qty{3.156e7}{s.year^{-1}}} = \qty{1.40e10}{years}$.
  The Universe is expanding.}
\significance{the first strong evidence that the Universe is expanding, not static. It
  points back to the Big Bang. Georges Lema\^{\i}tre predicted the expansion in 1927, so
  this is also called the Hubble--Lema\^{\i}tre law.}
\validity{
  \textbf{Validity:} it assumes redshift comes from the expansion. Use distant
        galaxies, because nearby ones have their own motion. Distances need a correct
        Cepheid scale: Hubble's value (about \qty{500}{km.s^{-1}.Mpc^{-1}}) was far too big
        because his distances were too small.\par
  \textbf{Reliability:} use many galaxies and data from different teams and methods.\par
  \textbf{Accuracy:} distance is the biggest uncertainty; a line of best fit
        through many points reduces random errors.
}
\end{experiment}

\begin{experiment}{Observing continuous, emission and absorption spectra}
\aim{To see and compare the three kinds of spectrum.}
\apparatus{hand-held spectroscope (or diffraction grating); a bright white filament lamp
  (a halogen or projector lamp); gas discharge tubes (hydrogen, helium, neon) with a
  high-voltage supply; a sodium flame (salt in a Bunsen flame). The flame is hot, but
  cooler than the lamp's filament, so its sodium gas absorbs light.
  \begin{figure}
    \centering
    \begin{tikzpicture}[x=1cm, y=1cm]
      \draw[hscink, line width=0.8pt] (0,0) circle (0.45);
      \draw[hscacc, line width=0.9pt, decorate, decoration={zigzag, amplitude=0.6mm,
        segment length=1.2mm}] (-0.2,0) -- (0.2,0);
      \node[diagramlabel, below, align=center, font=\footnotesize] at (0,-0.55)
        {filament lamp\\(or gas tube)};
      \fill[hscray!60] (2.7,-0.45) .. controls (2.35,0.05) and (2.6,0.35) .. (2.7,0.6)
        .. controls (2.8,0.35) and (3.05,0.05) .. (2.7,-0.45);
      \draw[hscmuted, line width=0.8pt] (2.55,-0.45) rectangle (2.85,-0.95);
      \node[diagramlabel, below, align=center, font=\footnotesize, text width=30mm]
        at (2.7,-1.0) {sodium flame (cooler than the filament): only for absorption};
      \draw[ray] (0.5,0) -- (4.8,0);
      \draw[object] (4.8,-0.32) rectangle (7.9,0.32);
      \draw[hscink, line width=1.4pt] (4.8,-0.12) -- (4.8,-0.32) (4.8,0.12) -- (4.8,0.32);
      \node[diagramlabel, above, font=\footnotesize] at (4.8,0.35) {slit};
      \draw[hscfield, line width=1pt] (7.6,-0.3) -- (7.6,0.3);
      \node[diagramlabel, above, font=\footnotesize] at (7.6,0.35) {grating};
      \node[diagramlabel, below, font=\footnotesize] at (6.35,-0.35) {spectroscope};
      \draw[hscink, line width=0.8pt] (9.05,0.3) .. controls (8.65,0.15) and (8.5,0) ..
        (8.5,0) .. controls (8.5,0) and (8.65,-0.15) .. (9.05,-0.3);
      \fill[hscink] (8.7,0) circle (0.07);
      \node[diagramlabel, right, font=\footnotesize] at (9.1,0) {eye};
    \end{tikzpicture}
    \caption{Viewing spectra. Remove the flame to see the lamp's continuous spectrum.}
    \label{ch1:fig:spectroscope}
  \end{figure}}
\method{\begin{enumerate}
  \item Look at the filament lamp through the spectroscope. Sketch the spectrum.
  \item Look at each gas tube. Record the colour and scale reading of each bright line.
  \item Put the sodium flame between the lamp and the spectroscope. Record any dark lines.
        Compare all readings with a table of known wavelengths.
\end{enumerate}}
\results{the lamp gives a continuous rainbow. Each gas gives its own bright lines
  (hydrogen: \qty{656}{nm} red, \qty{486}{nm} blue-green, 434 and \qty{410}{nm} violet).
  Sodium gives a dark line at \qty{589}{nm}, where its bright yellow line is. This dark
  line is faint; it is clearer with a brighter lamp. Another way to see dark lines: look
  at daylight reflected from a white cloud (never at the Sun). Its spectrum shows the
  Sun's own dark absorption lines (Fraunhofer lines).}
\conclusion{hot dense sources give continuous spectra, hot thin gases give bright lines,
  and cool gas in front of a hot source gives dark lines at the same wavelengths.}
\significance{this is how astronomers find what stars are made of and how hot they are.}
\validity{
  \textbf{Validity:} same spectroscope and viewing distance; block stray room light.\par
  \textbf{Reliability:} repeat each reading and compare with other groups.\par
  \textbf{Accuracy:} a hand spectroscope reads to only about \qty{\pm 5}{nm}; a
        grating on a spectrometer ($d\sin\theta = m\lambda$) is more accurate.
        \emph{Safety:} gas tubes use high voltage; never look at the Sun.
}
\end{experiment}

\begin{experiment}{Plotting an HR diagram from star data}
\aim{To plot stars by surface temperature and absolute magnitude, and to find the main
  groups of stars.}
\apparatus{a data table of stars (\cref{ch1:tab:hrdata}; approximate values from a star
  catalogue), and log graph paper or a spreadsheet.
  \begin{table}
    \caption{Data for ten real stars.}\label{ch1:tab:hrdata}
    \centering\small
    \begin{tabular}{@{}lrr@{\hspace{10mm}}lrr@{}}
      \toprule
      \colhead{Star} & \colhead{$T$ (K)} & \colhead{$M$} &
        \colhead{Star} & \colhead{$T$ (K)} & \colhead{$M$}\\
      \midrule
      Sun & \num{5772} & $+4.8$ & Betelgeuse & \num{3600} & $-5.9$\\
      Sirius A & \num{9940} & $+1.4$ & Aldebaran & \num{3900} & $-0.6$\\
      Sirius B & \num{25000} & $+11.2$ & Procyon B & \num{7700} & $+13.0$\\
      Vega & \num{9600} & $+0.6$ & Barnard's star & \num{3130} & $+13.2$\\
      Rigel & \num{12100} & $-7.8$ & Proxima Centauri & \num{3040} & $+15.5$\\
      \bottomrule
    \end{tabular}
  \end{table}}
\method{\begin{enumerate}
  \item Draw the temperature axis as a log scale, with hot on the \emph{left} (for
        example from \qty{30000}{K} down to \qty{2500}{K}).
  \item Draw the absolute magnitude axis with the negative (bright) numbers at the top.
  \item Plot each star and write its name next to it.
  \item Shade the diagonal band of the main sequence. Circle the giants, the supergiants
        and the white dwarfs.
\end{enumerate}}
\results{\begin{figure}
    \centering
    \includegraphics{ch1_hrplot.pdf}
    \caption{The ten stars of \cref{ch1:tab:hrdata} on an HR
      diagram.}\label{ch1:fig:hrplot}
  \end{figure}}
\conclusion{the Sun, Sirius A, Vega, Barnard's star and Proxima Centauri lie on a
  diagonal band: the main sequence. Aldebaran (a giant) and Betelgeuse and Rigel
  (supergiants) lie above it. They are very luminous for their temperature, so they are
  large. Sirius B and Procyon B lie below and to the left. They are hot but dim, so they
  are small white dwarfs.}
\significance{Hertzsprung and Russell found these groups in the same way (about
  1911--1913). The groups showed that stars change as they age, and that the main
  sequence is the long, steady stage of a star's life.}
\validity{
  \textbf{Validity:} use absolute magnitude, not apparent magnitude, so that distance does
        not change a star's place. (Stars from one cluster are all at about the same
        distance, so apparent magnitude works for a cluster.)\par
  \textbf{Reliability:} use many stars (a large sample) from a trusted catalogue.\par
  \textbf{Accuracy:} an error in a star's distance gives an error in its absolute
        magnitude. Temperatures from spectra are good to a few per cent.
}
\end{experiment}

